\chapter[Divisors from Fans]
{Divisors from Fans:\\ Rays, Characters, and\\ Piecewise-Linear Functions}
\label{chap:divisors-from-fans}

\begin{chapteroverview}
The orbit--cone correspondence tells us that a ray
\[
\rho\in\Sigma(1)
\]
corresponds to a codimension-one torus orbit. Its closure is therefore
an irreducible, torus-invariant prime divisor
\[
D_\rho=\overline{O(\rho)}.
\]
This is the first principle of the chapter:
\[
\boxed{\text{Rays record the irreducible components of the toric boundary.}}
\]

We will use the primitive ray generators \(u_\rho\) to measure zeros
and poles of characters:
\[
\operatorname{div}(\character{m})
=
\sum_{\rho\in\Sigma(1)}
\langle m,u_\rho\rangle D_\rho.
\]
These principal divisors provide the relations among the invariant
prime divisors and lead to a combinatorial description of the divisor
class group.

Cartier divisors will then be described by one local character on
each cone. Compatibility of those local equations is equivalent to
piecewise linearity on the fan. Thus the chapter develops the passage
\[
\begin{array}{c}
\text{rays and integral linear functions}
\\[-1mm]
\downarrow
\\[-1mm]
\text{prime divisors and local equations}
\\[-1mm]
\downarrow
\\[-1mm]
\text{divisor classes and support functions}.
\end{array}
\]

\end{chapteroverview}

\begin{threedictionaries}
Divisors add codimension-one information to the three dictionaries.

\medskip

\begin{center}
\small
\begin{tabularx}{\linewidth}{@{}
  >{\raggedright\arraybackslash}X
  >{\raggedright\arraybackslash}X
  >{\raggedright\arraybackslash}X@{}}
\toprule
\textbf{Algebra}
&
\textbf{Geometry}
&
\textbf{Combinatorics}
\\
\midrule
a valuation along \(D_\rho\)
&
order of zero or pole
&
the pairing with \(u_\rho\)
\\[1mm]
a principal divisor
&
relations among boundary components
&
an element \(m\in M\)
\\[1mm]
local monomial equations
&
a Cartier divisor
&
linear functions on cones
\\[1mm]
compatible local equations
&
a global divisor class
&
an integral piecewise-linear function
\\
\bottomrule
\end{tabularx}
\end{center}
\end{threedictionaries}

\begin{surfacelaboratory}
For a toric surface, every ray gives a torus-invariant curve and every
two-dimensional cone gives a torus-fixed point.

Throughout the chapter, we will repeatedly ask:
\begin{itemize}
\item Which boundary curve corresponds to each ray?
\item Which curves meet?
\item Which integer relations among the ray generators become linear
equivalences among divisors?
\item Which divisors admit compatible local equations?
\end{itemize}
\end{surfacelaboratory}

\section{Rays become prime divisors}
\label{sec:rays-become-prime-divisors}
\index{ray!invariant divisor}
\index{invariant prime divisor}
\index{toric boundary}
\index{boundary incidence}

The previous chapter indexed every torus orbit by a cone:
\[
\sigma\in\Sigma
\quad\longleftrightarrow\quad
O(\sigma)\subseteq X_\Sigma,
\]
with
\[
\dim O(\sigma)
=
\rank N-\dim\sigma.
\]
We now focus on the cones of dimension one.

\subsection{A brief divisor dictionary}

\begin{agcheatsheet}
Let \(X\) be an irreducible normal variety.

A \emph{prime divisor} is an irreducible subvariety
\[
D\subseteq X
\]
of codimension one.

A \emph{Weil divisor} is a finite formal sum
\[
D=\sum_i a_iD_i,
\qquad
a_i\in\ZZ,
\]
of prime divisors. The group of Weil divisors is denoted
\[
\Div(X).
\]

The divisor is \emph{effective}, written \(D\geq0\), if every
\(a_i\geq0\). Its support is
\[
\Supp(D)
=
\bigcup_{a_i\neq0}D_i.
\]

At this stage, a Weil divisor should be viewed as codimension-one
geometry with integer multiplicities. Zeros, poles, and linear
equivalence will enter in the next section.
\end{agcheatsheet}

Normality is important because codimension-one subvarieties carry
well-behaved orders of vanishing. Toric varieties constructed from
fans are normal, so Weil divisors are available throughout this
chapter.

\subsection{One divisor for every ray}

Let
\[
\rho\in\Sigma(1).
\]
The orbit--cone correspondence gives
\[
\dim O(\rho)
=
\rank N-1.
\]
Its closure is therefore irreducible and has codimension one.

\begin{definition}
\label{def:torus-invariant-prime-divisor}
The \emph{torus-invariant prime divisor associated to \(\rho\)} is
\[
D_\rho
=
V(\rho)
=
\overline{O(\rho)}
\subseteq X_\Sigma.
\]
\end{definition}

The notation \(V(\rho)\) emphasizes orbit closure; the notation
\(D_\rho\) emphasizes codimension one. We will use \(D_\rho\) from now
on.

\begin{theorem}[Rays and invariant prime divisors]
\label{thm:rays-invariant-prime-divisors}
The assignment
\[
\rho\longmapsto D_\rho
\]
is a bijection
\[
\Sigma(1)
\longleftrightarrow
\left\{
\begin{array}{c}
\text{torus-invariant}\\
\text{prime divisors of }X_\Sigma
\end{array}
\right\}.
\]
\end{theorem}

\noindent
\textit{Reference.}
The notation \(D_\rho=\overline{O(\rho)}\) is introduced at the
beginning of Cox--Little--Schenck, \S4.1. The characterization of all
torus-invariant Weil divisors is Exercise~4.1.1 there.

\begin{proof}
Every \(D_\rho\) is irreducible because it is the closure of one
orbit, and it has codimension one by the orbit-dimension formula. It
is torus-invariant because the closure of a torus orbit is preserved
by the torus action.

Conversely, let \(D\subseteq X_\Sigma\) be a torus-invariant prime
divisor. Since the torus acts transitively on itself, the intersection
\[
D\cap T
\]
is either empty or all of \(T\). It cannot be all of \(T\), because
\(D\) is a proper closed subset. Hence
\[
D\subseteq X_\Sigma\setminus T.
\]

Every nonzero cone contains a ray. If
\[
\sigma\neq\{0\}
\qquad\text{and}\qquad
\rho\preceq\sigma,
\]
then the orbit--cone correspondence gives
\[
O(\sigma)\subseteq\overline{O(\rho)}=D_\rho.
\]
Therefore
\[
X_\Sigma\setminus T
=
\bigcup_{\rho\in\Sigma(1)}D_\rho.
\]
The irreducible divisor \(D\) must lie in one of these finitely many
closed sets. Since \(D\) and \(D_\rho\) are both irreducible of
codimension one, they are equal.
\end{proof}

Thus the fan lists every invariant codimension-one subvariety before
we write a single equation.

\subsection{The toric boundary}

\begin{definition}
The \emph{toric boundary} is the complement of the dense torus:
\[
\partial X_\Sigma
=
X_\Sigma\setminus T.
\]
As a reduced Weil divisor, it is
\[
\partial X_\Sigma
=
\sum_{\rho\in\Sigma(1)}D_\rho.
\]
\end{definition}

The same symbol is traditionally used for the closed subset and for
its associated reduced divisor; the intended meaning will be clear
from context.

\begin{example}[The boundary of \(\AAA^2\)]
Let
\[
\AAA^2=\Spec\Bbbk[x,y],
\]
and let the rays of its first-quadrant fan be
\[
\rho_1=\Cone{e_1},
\qquad
\rho_2=\Cone{e_2}.
\]
The coordinate axes are precisely the corresponding invariant prime
divisors:
\[
V(x)=D_{\rho_1},
\qquad
V(y)=D_{\rho_2}.
\]
Therefore, as a closed subset,
\[
\partial\AAA^2
=
V(x)\cup V(y),
\]
while as a reduced Weil divisor,
\[
\partial\AAA^2
=
D_{\rho_1}+D_{\rho_2}
=
V(x)+V(y).
\]
\end{example}

For a complete toric variety, every direction has been filled in by
the fan, but the toric boundary remains: completeness does not mean
that the dense torus has no boundary. It means that every
one-parameter direction approaches some boundary orbit.

\begin{pauseandthink}
The fan of \(\PP^1\) has two rays. Identify the corresponding prime
divisors, remembering that prime divisors on a curve are points.

What is the toric boundary of \(\PP^1\)?
\end{pauseandthink}

\subsection{The fan remembers incidence}

The fan records not only the boundary components but also how they
meet.

\begin{proposition}[Boundary incidence]
\label{prop:boundary-incidence}
Let \(\rho_1,\ldots,\rho_s\in\Sigma(1)\). Then
\[
D_{\rho_1}\cap\cdots\cap D_{\rho_s}\neq\varnothing
\]
if and only if the rays \(\rho_1,\ldots,\rho_s\) are all contained in
a cone of \(\Sigma\).

For a toric surface, two distinct boundary curves meet if and only if
their rays are adjacent. When they meet, their intersection is the
torus-fixed point corresponding to the two-dimensional cone that they
span.
\end{proposition}

\begin{proof}
Suppose the rays lie in a cone \(\sigma\). Then
\[
\rho_i\preceq\sigma
\]
for every \(i\), so
\[
O(\sigma)\subseteq
\overline{O(\rho_i)}
=
D_{\rho_i}.
\]
Thus the intersection is nonempty.

Conversely, let \(p\) lie in the intersection. The point \(p\) lies
in a unique orbit \(O(\sigma)\). Since
\[
O(\sigma)\subseteq D_{\rho_i},
\]
the orbit-closure relation implies
\[
\rho_i\preceq\sigma
\]
for every \(i\).

On a surface, adjacent rays span a two-dimensional cone, whose orbit
is one point. No larger cone exists, so this fixed point is the whole
intersection.
\end{proof}

This proposition turns the cyclic ordering of rays in a surface fan
into the incidence graph of its boundary curves.

\subsection{The coordinate triangle in \texorpdfstring{\(\PP^2\)}{P2}}

Let the rays of the fan of \(\PP^2\) be
\[
\rho_0=\Cone{-e_1-e_2},
\qquad
\rho_1=\Cone{e_1},
\qquad
\rho_2=\Cone{e_2}.
\]
The associated invariant divisors are the coordinate lines
\[
D_{\rho_i}
=
\{z_i=0\}.
\]

\begin{figure}[ht]
\centering
\begin{tikzpicture}[scale=1.0]
  % fan
  \begin{scope}[xshift=-4.4cm]
    \fill[DictionaryBlue!9]
      (0,0)--(2.3,0)--(0,2.3)--cycle;
    \fill[SurfaceGreen!9]
      (0,0)--(0,2.3)--(-1.7,-1.7)--cycle;
    \fill[ExampleGold!12]
      (0,0)--(-1.7,-1.7)--(2.3,0)--cycle;

    \draw[very thick,DictionaryBlue,-{Latex[length=2.3mm]}]
      (0,0)--(2.65,0);
    \draw[very thick,SurfaceGreen,-{Latex[length=2.3mm]}]
      (0,0)--(0,2.65);
    \draw[very thick,ExampleGold,-{Latex[length=2.3mm]}]
      (0,0)--(-2.0,-2.0);

    \node[DictionaryBlue,below] at (2.45,0) {\(\rho_1\)};
    \node[SurfaceGreen,left] at (0,2.45) {\(\rho_2\)};
    \node[ExampleGold,above left] at (-1.82,-1.82) {\(\rho_0\)};
    \node[font=\bfseries] at (0.2,3.25) {rays of the fan};
  \end{scope}

  \node[font=\Large] at (0,0.4) {\(\longleftrightarrow\)};

  % coordinate triangle
  \begin{scope}[xshift=4.3cm]
    \coordinate (p0) at (0,2.5);
    \coordinate (p1) at (-2.35,-1.5);
    \coordinate (p2) at (2.35,-1.5);

    \draw[very thick,ExampleGold] (p1)--(p2);
    \draw[very thick,DictionaryBlue] (p0)--(p2);
    \draw[very thick,SurfaceGreen] (p0)--(p1);

    \fill (p0) circle (2.2pt);
    \fill (p1) circle (2.2pt);
    \fill (p2) circle (2.2pt);

    \node[above] at (p0) {\([1:0:0]\)};
    \node[below left] at (p1) {\([0:1:0]\)};
    \node[below right] at (p2) {\([0:0:1]\)};

    \node[ExampleGold,below] at (0,-1.5)
      {\(D_{\rho_0}=\{z_0=0\}\)};
    \node[DictionaryBlue,rotate=-60] at (1.62,0.62)
      {\(D_{\rho_1}=\{z_1=0\}\)};
    \node[SurfaceGreen,rotate=60] at (-1.62,0.62)
      {\(D_{\rho_2}=\{z_2=0\}\)};

    \node[font=\bfseries] at (0,3.25)
      {toric boundary of \(\PP^2\)};
  \end{scope}
\end{tikzpicture}
\caption{The three rays of the fan correspond to the three coordinate
lines in \(\PP^2\). Adjacent rays span maximal cones, just as pairs of
coordinate lines meet at coordinate points. The triangle on the right
is an incidence picture of the boundary.}
\label{fig:p2-rays-coordinate-divisors}
\end{figure}

Thus
\[
\partial\PP^2
=
D_{\rho_0}+D_{\rho_1}+D_{\rho_2}.
\]
Every pair of rays spans a maximal cone, so every pair of coordinate
lines meets. The three lines have no common point, matching the fact
that no cone contains all three rays.

\begin{surfacelaboratory}
Read the boundary of \(\PP^2\) from
\cref{fig:p2-rays-coordinate-divisors}.

\begin{enumerate}[label=(\alph*),wide=0pt]
\item Match each maximal cone with the coordinate point at which its
two boundary divisors meet.

\item Explain why
\[
D_{\rho_0}\cap D_{\rho_1}\cap D_{\rho_2}
=
\varnothing
\]
without using homogeneous coordinates.

\item Remove the ray \(\rho_0\) and every cone containing it. Which
toric open subset remains, and which coordinate boundary components
does it contain?
\end{enumerate}
\end{surfacelaboratory}

\subsection{A new ray creates the exceptional divisor}

Consider again the fan of the blowup
\[
\pi:\Blp{\AAA^2}{0}\longrightarrow\AAA^2.
\]
It is obtained by inserting
\[
\rho_E=\Cone{e_1+e_2}
\]
between
\[
\rho_1=\Cone{e_1}
\qquad\text{and}\qquad
\rho_2=\Cone{e_2}.
\]

\begin{figure}[ht]
\centering
\begin{tikzpicture}[scale=1.0]
  % fan
  \begin{scope}[xshift=-4.5cm]
    \fill[DictionaryBlue!10]
      (0,0)--(2.6,0)--(1.3,1.3)--cycle;
    \fill[SurfaceGreen!10]
      (0,0)--(1.3,1.3)--(0,2.6)--cycle;

    \draw[very thick,DictionaryBlue,-{Latex[length=2.3mm]}]
      (0,0)--(2.9,0);
    \draw[very thick,PredictionTeal,-{Latex[length=2.3mm]}]
      (0,0)--(2.25,2.25);
    \draw[very thick,SurfaceGreen,-{Latex[length=2.3mm]}]
      (0,0)--(0,2.9);

    \node[DictionaryBlue,below] at (2.7,0) {\(\rho_1\)};
    \node[PredictionTeal,above right] at (2.0,2.0) {\(\rho_E\)};
    \node[SurfaceGreen,left] at (0,2.7) {\(\rho_2\)};
    \node[font=\bfseries] at (1.0,3.4) {blowup fan};
  \end{scope}

  \node[font=\Large] at (0,0.5) {\(\longleftrightarrow\)};

  % boundary incidence
  \begin{scope}[xshift=4.4cm]
    \draw[very thick,DictionaryBlue] (-2.6,-0.2)--(-0.65,-0.2);
    \draw[very thick,PredictionTeal] (-0.65,-0.2) arc[start angle=180,end angle=0,
      x radius=1.3,y radius=1.3];
    \draw[very thick,SurfaceGreen] (1.95,-0.2)--(3.4,-0.2);

    \fill (-0.65,-0.2) circle (2.2pt);
    \fill (1.95,-0.2) circle (2.2pt);

    \node[DictionaryBlue,below] at (-1.75,-0.2)
      {\(D_{\rho_1}\)};
    \node[PredictionTeal,above] at (0.65,1.15)
      {\(E=D_{\rho_E}\simeq\PP^1\)};
    \node[SurfaceGreen,below] at (2.7,-0.2)
      {\(D_{\rho_2}\)};

    \node[align=center] at (0.65,-1.5)
      {the two strict transforms no longer meet;\\
       each meets the exceptional divisor};
    \node[font=\bfseries] at (0.45,3.4)
      {boundary incidence};
  \end{scope}
\end{tikzpicture}
\caption{The new ray of the blowup fan creates the exceptional
divisor. The subdivision also changes incidence: the strict transforms
of the coordinate axes are separated by \(E\).}
\label{fig:blowup-rays-boundary-divisors}
\end{figure}

The old rays correspond to the strict transforms of the coordinate
axes. The new ray corresponds to the exceptional curve
\[
E=D_{\rho_E}.
\]
Because \(\rho_E\) lies in two maximal cones, its star is the complete
one-dimensional fan, so
\[
E\simeq\PP^1.
\]

The original rays \(\rho_1,\rho_2\) no longer span a cone. Hence
\[
D_{\rho_1}\cap D_{\rho_2}=\varnothing.
\]
Instead,
\[
D_{\rho_1}\cap E\neq\varnothing,
\qquad
E\cap D_{\rho_2}\neq\varnothing.
\]

\begin{pauseandthink}
Compare the fans before and after the blowup.

Which single combinatorial change simultaneously explains:
\begin{enumerate}[label=(\alph*),wide=0pt]
\item the appearance of a new boundary curve;
\item the appearance of one additional torus-fixed point;
\item the separation of the strict transforms of the coordinate axes?
\end{enumerate}
\end{pauseandthink}

The fan tells us which prime divisors exist, but not yet how to
write equations for them or compare them.

\begin{guidingquestion}
How can the integer \(\langle m,u_\rho\rangle\) measure the order of
vanishing or pole of \(\character{m}\) along \(D_\rho\)?
\end{guidingquestion}

\section{Characters and principal divisors}
\label{sec:characters-principal-divisors}
\index{principal divisor}
\index{divisor of a character}
\index{order of vanishing}
\index{linear equivalence}

A character
\[
\character{m}:T\longrightarrow\Gm,
\qquad
m\in M,
\]
is regular and nowhere zero on the dense torus. Since
\[
T\subseteq X_\Sigma
\]
is dense, \(\character{m}\) is a rational function on \(X_\Sigma\).

It can acquire zeros and poles only when we approach the toric
boundary. The primitive ray generators tell us exactly what happens.

\subsection{The divisor of a character}

\begin{agcheatsheet}
Let \(X\) be a normal irreducible variety, let \(D\subseteq X\) be a
prime divisor, and let
\[
f\in\Bbbk(X)^\times
\]
be a nonzero rational function.

The local ring at the generic point of \(D\) is a discrete valuation
ring. Consequently, near the generic point of \(D\), one can write
\[
f=\pi^\ell u,
\]
where \(\pi\) is a local equation for \(D\), \(u\) is a unit, and
\(\ell\in\ZZ\). We define
\[
\ord_D(f)=\ell.
\]

Thus
\[
\begin{array}{c|c}
\ord_D(f)>0 & f\text{ has a zero along }D,\\
\ord_D(f)=0 & f\text{ is a unit at the generic point of }D,\\
\ord_D(f)<0 & f\text{ has a pole along }D.
\end{array}
\]

The \emph{principal divisor} of \(f\) is
\[
\operatorname{div}(f)
=
\sum_D\ord_D(f)D,
\]
where the sum runs over all prime divisors of \(X\). Only finitely
many coefficients are nonzero.
\end{agcheatsheet}

The order is measured at the generic point of the divisor. A function
may vanish at isolated points of \(D\) without vanishing
\emph{along} \(D\).

For a toric divisor \(D_\rho\), we abbreviate
\[
\ord_{D_\rho}(f)
\]
by
\[
\ord_\rho(f).
\]

\begin{proposition}[The order of a character]
\label{prop:order-character-ray}
Let \(\rho\in\Sigma(1)\) have primitive generator \(u_\rho\). Then,
for every \(m\in M\),
\[
\ord_{D_\rho}(\character{m})
=
\langle m,u_\rho\rangle.
\]
\end{proposition}

\noindent
\textit{Reference.}
This is Cox--Little--Schenck, Proposition~4.1.1.

\begin{proof}[Conceptual proof sketch]
Because \(u_\rho\) is primitive, it can be completed to a basis of
\(N\). In the resulting toric coordinates,
\[
U_\rho\simeq\AAA^1\times\Gm^{n-1},
\qquad
D_\rho\cap U_\rho=V(x_1).
\]
The character \(\character{m}\) is then \(x_1^{\langle
m,u_\rho\rangle}\) times a unit along \(D_\rho\). Its order is
therefore \(\langle m,u_\rho\rangle\).
\end{proof}

\begin{remark}[Why the primitive generator matters]
The valuation \(\ord_{D_\rho}\) is normalized so that a local
parameter has order one. If we replaced the primitive generator by
\[
v=k u_\rho,
\qquad k>1,
\]
then
\[
\langle m,v\rangle
=
k\langle m,u_\rho\rangle,
\]
which would multiply every order by \(k\). The ray determines the
divisor, but its \emph{primitive} lattice generator gives the correctly
normalized valuation.
\end{remark}

The formula has a useful one-parameter interpretation. Along
\(\oneps{u_\rho}\),
\[
\character{m}(\oneps{u_\rho}(t))
=
t^{\langle m,u_\rho\rangle}.
\]
The same integer therefore controls both the limit along the ray and
the order along the corresponding divisor.

\begin{theorem}[Principal divisor formula]
\label{thm:principal-divisor-character}
For every \(m\in M\),
\[
\boxed{
\operatorname{div}(\character{m})
=
\sum_{\rho\in\Sigma(1)}
\langle m,u_\rho\rangle D_\rho.
}
\]
\end{theorem}

\noindent
\textit{Reference.}
This is Cox--Little--Schenck, Proposition~4.1.2.

\begin{pauseandthink}
Fill in the proof. There are only two steps:
\begin{enumerate}[label=(\alph*),wide=0pt]
\item Explain why \(\operatorname{div}(\character{m})\) is supported
on the toric boundary.
\item Use \cref{prop:order-character-ray} to compute the coefficient
of every \(D_\rho\).
\end{enumerate}
\end{pauseandthink}

The signs have a direct geometric meaning:
\[
\begin{array}{c@{\qquad}c}
\langle m,u_\rho\rangle>0
&
D_\rho\text{ occurs in the zero divisor},
\\[1mm]
\langle m,u_\rho\rangle=0
&
\character{m}\text{ is generically nonzero along }D_\rho,
\\[1mm]
\langle m,u_\rho\rangle<0
&
D_\rho\text{ occurs in the pole divisor}.
\end{array}
\]

Thus a linear function on \(N_\RR\) becomes a signed collection of
multiplicities on the boundary.

\subsection{First examples}

\begin{example}[Monomials on \(\AAA^2\)]
Let
\[
\rho_1=\Cone{e_1},
\qquad
\rho_2=\Cone{e_2},
\]
so that
\[
D_{\rho_1}=V(x),
\qquad
D_{\rho_2}=V(y).
\]
For
\[
m=am_1+bm_2,
\]
we have
\[
\character{m}=x^ay^b
\]
as a rational function, and
\[
\operatorname{div}(x^ay^b)
=
aD_{\rho_1}+bD_{\rho_2}.
\]

In particular,
\[
\operatorname{div}(x)=D_{\rho_1},
\qquad
\operatorname{div}(y)=D_{\rho_2},
\]
while
\[
\operatorname{div}(x/y)
=
D_{\rho_1}-D_{\rho_2}.
\]
\end{example}

\begin{example}[Zeros and poles on \(\PP^1\)]
Let
\[
\rho_+=\RR_{\geq0},
\qquad
\rho_-=\RR_{\leq0},
\]
with primitive generators \(1\) and \(-1\). The character
\[
t=\character{1}
\]
satisfies
\[
\operatorname{div}(t)
=
D_{\rho_+}-D_{\rho_-}.
\]
Under the usual identification,
\[
D_{\rho_+}=\{0\},
\qquad
D_{\rho_-}=\{\infty\},
\]
so this is the familiar formula
\[
\operatorname{div}(t)
=
[0]-[\infty].
\]
\end{example}

\subsection{The coordinate characters on \texorpdfstring{\(\PP^2\)}{P2}}

Let
\[
u_{\rho_0}=-e_1-e_2,
\qquad
u_{\rho_1}=e_1,
\qquad
u_{\rho_2}=e_2.
\]
Write
\[
x=\character{m_1}=\frac{z_1}{z_0},
\qquad
y=\character{m_2}=\frac{z_2}{z_0}.
\]

\begin{figure}[ht]
\centering
\begin{tikzpicture}[scale=0.9]
  % x
  \begin{scope}[xshift=-4.4cm]
    \fill[DictionaryBlue!8] (0,0)--(2.3,0)--(0,2.3)--cycle;
    \fill[LightGray] (0,0)--(0,2.3)--(-1.7,-1.7)--cycle;
    \fill[WarningRed!8] (0,0)--(-1.7,-1.7)--(2.3,0)--cycle;

    \draw[very thick,DictionaryBlue,-{Latex[length=2.3mm]}]
      (0,0)--(2.65,0);
    \draw[very thick,black!45,-{Latex[length=2.3mm]}]
      (0,0)--(0,2.65);
    \draw[very thick,WarningRed,-{Latex[length=2.3mm]}]
      (0,0)--(-2.0,-2.0);

    \node[DictionaryBlue,below] at (2.45,0)
      {\(\langle m_1,u_{\rho_1}\rangle=1\)};
    \node[black!60,left] at (0,2.45)
      {\(\langle m_1,u_{\rho_2}\rangle=0\)};
    \node[WarningRed,above left] at (-1.8,-1.8)
      {\(\langle m_1,u_{\rho_0}\rangle=-1\)};

    \node[font=\bfseries] at (0.2,3.25)
      {\(x=z_1/z_0\)};
    \node[align=center] at (0.15,-2.75)
      {\(\operatorname{div}(x)=D_{\rho_1}-D_{\rho_0}\)};
  \end{scope}

  % y
  \begin{scope}[xshift=4.4cm]
    \fill[LightGray] (0,0)--(2.3,0)--(0,2.3)--cycle;
    \fill[SurfaceGreen!8] (0,0)--(0,2.3)--(-1.7,-1.7)--cycle;
    \fill[WarningRed!8] (0,0)--(-1.7,-1.7)--(2.3,0)--cycle;

    \draw[very thick,black!45,-{Latex[length=2.3mm]}]
      (0,0)--(2.65,0);
    \draw[very thick,SurfaceGreen,-{Latex[length=2.3mm]}]
      (0,0)--(0,2.65);
    \draw[very thick,WarningRed,-{Latex[length=2.3mm]}]
      (0,0)--(-2.0,-2.0);

    \node[black!60,below] at (2.45,0)
      {\(\langle m_2,u_{\rho_1}\rangle=0\)};
    \node[SurfaceGreen,left] at (0,2.45)
      {\(\langle m_2,u_{\rho_2}\rangle=1\)};
    \node[WarningRed,above left] at (-1.8,-1.8)
      {\(\langle m_2,u_{\rho_0}\rangle=-1\)};

    \node[font=\bfseries] at (0.2,3.25)
      {\(y=z_2/z_0\)};
    \node[align=center] at (0.15,-2.75)
      {\(\operatorname{div}(y)=D_{\rho_2}-D_{\rho_0}\)};
  \end{scope}
\end{tikzpicture}
\caption{The coefficients in the divisor of a character are obtained
by evaluating the corresponding linear function on the primitive ray
generators. Positive coefficients record zeros, red negative
coefficients record poles, and gray zero coefficients contribute
nothing.}
\label{fig:p2-character-divisors}
\end{figure}

The figure gives
\[
\operatorname{div}(x)
=
D_{\rho_1}-D_{\rho_0},
\qquad
\operatorname{div}(y)
=
D_{\rho_2}-D_{\rho_0}.
\]
More generally,
\[
\operatorname{div}(x^ay^b)
=
aD_{\rho_1}
+
bD_{\rho_2}
-
(a+b)D_{\rho_0}.
\]

This agrees with homogeneous coordinates:
\[
x^ay^b
=
\frac{z_1^az_2^b}{z_0^{a+b}}.
\]
The numerator records the zeros, and the denominator records the
poles.

\subsection{Principal divisors and linear equivalence}

\begin{definition}
Two Weil divisors \(D,D'\) are \emph{linearly equivalent}, written
\[
D\sim D',
\]
if
\[
D-D'=\operatorname{div}(f)
\]
for some \(f\in\Bbbk(X)^\times\).

A divisor of the form \(\operatorname{div}(f)\) is called
\emph{principal}. Thus every principal divisor is linearly equivalent
to zero.
\end{definition}

The character calculations on \(\PP^2\) imply
\[
D_{\rho_1}\sim D_{\rho_0},
\qquad
D_{\rho_2}\sim D_{\rho_0}.
\]
Hence the three coordinate lines are linearly equivalent.

This is the first appearance of a general principle:
\[
\text{characters}
\quad\leadsto\quad
\text{relations among invariant divisors}.
\]

\begin{pauseandthink}
Let
\[
m=2m_1-m_2.
\]
Use \cref{fig:p2-character-divisors}, or pair directly with the three
ray generators, to compute
\[
\operatorname{div}(x^2/y).
\]
Which coordinate lines occur as zeros, and which occur as poles?
\end{pauseandthink}

\subsection{The blowup remembers total transforms}

Let
\[
\pi:\Blp{\AAA^2}{0}\longrightarrow\AAA^2
\]
be the toric blowup. Its rays have primitive generators
\[
e_1,\qquad
e_1+e_2,\qquad
e_2.
\]
Write
\[
D_1=D_{\Cone{e_1}},
\qquad
E=D_{\Cone{e_1+e_2}},
\qquad
D_2=D_{\Cone{e_2}}.
\]

The pairing table is
\[
\begin{array}{c|ccc}
&
e_1
&
e_1+e_2
&
e_2
\\
\hline
m_1
&
1&1&0
\\
m_2
&
0&1&1.
\end{array}
\]
Therefore
\[
\operatorname{div}(x)
=
D_1+E,
\qquad
\operatorname{div}(y)
=
D_2+E.
\]

\begin{surfacelaboratory}
Interpret these formulas geometrically.

\begin{enumerate}[label=(\alph*),wide=0pt]
\item Explain why \(D_1\) and \(D_2\) are the strict transforms of the
coordinate axes.

\item Explain why \(E\) appears in both
\[
\operatorname{div}(x)
\qquad\text{and}\qquad
\operatorname{div}(y).
\]
Relate this to the fact that both coordinate axes pass through the
blown-up point.

\item Compute
\[
\operatorname{div}(x/y).
\]
Why does the exceptional divisor cancel?

\item Use the formulas to obtain the linear equivalences
\[
D_1\sim -E,
\qquad
D_2\sim -E
\]
on the noncomplete surface \(\Blp{\AAA^2}{0}\).
\end{enumerate}
\end{surfacelaboratory}

\begin{warningbox}
The equality
\[
\operatorname{div}(x)=D_1+E
\]
is an equality of divisors on the blowup. It should not be confused
with the zero set of the function \(x\) on the original affine plane.

The divisor on the blowup records the total transform: the strict
transform \(D_1\) and the exceptional contribution \(E\).
\end{warningbox}

We now know the torus-invariant prime divisors and the relations among
them produced by rational functions. The next step is to decide when
two divisors should count as representing the same geometry.

\begin{guidingquestion}
Why might it be useful to regard two divisors as equivalent whenever
their difference is the divisor of a rational function?

What geometric information would you expect the quotient
\[
\Div(X_\Sigma)/
\Prin(X_\Sigma)
\]
to retain, and what information would it deliberately forget?
\end{guidingquestion}


\section{Weil divisors and the divisor class group}
\label{sec:weil-divisors-class-group}
\index{Weil divisor}
\index{divisor class group}
\index{class-group sequence}
\index{torsion!in the class group}

The preceding sections produced two pieces of information:

\[
\rho\in\Sigma(1)
\quad\longmapsto\quad
D_\rho,
\]
and
\[
m\in M
\quad\longmapsto\quad
\operatorname{div}(\character{m})
=
\sum_{\rho\in\Sigma(1)}
\langle m,u_\rho\rangle D_\rho.
\]

We now organize the invariant prime divisors modulo the relations
coming from characters.

\subsection{Invariant Weil divisors}

\begin{definition}
The group of \emph{torus-invariant Weil divisors} is
\[
\Div_T(X_\Sigma)
=
\bigoplus_{\rho\in\Sigma(1)}\ZZ D_\rho.
\]
Thus an invariant Weil divisor has the form
\[
D=\sum_{\rho\in\Sigma(1)}a_\rho D_\rho,
\qquad
a_\rho\in\ZZ.
\]
\end{definition}

Since the divisors \(D_\rho\) are distinct prime divisors, they form a
basis of this free abelian group. After ordering the rays, we may
identify
\[
\Div_T(X_\Sigma)
\simeq
\ZZ^{\Sigma(1)}.
\]

\begin{definition}
The \emph{divisor class group} is
\[
\Cl(X)
=
\Div(X)/
\Prin(X),
\]
the group of Weil divisors modulo principal divisors.

The class of a divisor \(D\) is denoted
\[
[D]\in\Cl(X).
\]
Thus
\[
[D]=[D']
\quad\Longleftrightarrow\quad
D\sim D'.
\]
\end{definition}

At first sight, computing \(\Cl(X_\Sigma)\) seems to
require all prime divisors of \(X_\Sigma\), not merely the invariant
ones. The torus action makes the answer much simpler.

\begin{theorem}[The class-group sequence]
\label{thm:toric-class-group-sequence}
There is an exact sequence
\[
M
\xrightarrow{\;\operatorname{div}\;}
\Div_T(X_\Sigma)
\longrightarrow
\Cl(X_\Sigma)
\longrightarrow
0,
\]
where
\[
m\longmapsto
\operatorname{div}(\character{m})
=
\sum_{\rho\in\Sigma(1)}
\langle m,u_\rho\rangle D_\rho,
\]
and the second map sends an invariant divisor to its divisor class.

The first map is injective if and only if the primitive ray
generators span \(N_\RR\). Equivalently, the sequence is short exact
precisely when \(X_\Sigma\) has no torus factor.
\end{theorem}

\noindent
\textit{Reference.}
This is Cox--Little--Schenck, Theorem~4.1.3.

\begin{proof}[Conceptual proof sketch]
The complement of the dense torus is the union of the invariant prime
divisors:
\[
X_\Sigma\setminus T
=
\bigcup_{\rho\in\Sigma(1)}D_\rho.
\]
Since the coordinate ring
\[
\Bbbk[M]
\]
of the torus is a Laurent polynomial ring, it is a unique
factorization domain and
\[
\Cl(T)=0.
\]
The localization sequence for divisor class groups therefore says
that every divisor class on \(X_\Sigma\) has a representative
supported on the boundary. This proves surjectivity:
\[
\Div_T(X_\Sigma)
\twoheadrightarrow
\Cl(X_\Sigma).
\]

Now suppose an invariant divisor is principal:
\[
D=\operatorname{div}(f).
\]
Its support misses the torus, so the restriction of \(f\) to \(T\)
has neither zeros nor poles. The units of \(\Bbbk[M]\) are
\[
c\character{m},
\qquad
c\in\Bbbk^\times,\quad m\in M.
\]
Hence
\[
D=\operatorname{div}(\character{m}),
\]
which identifies the kernel.

Finally,
\[
\operatorname{div}(\character{m})=0
\]
exactly when
\[
\langle m,u_\rho\rangle=0
\qquad
\text{for every }\rho.
\]
Thus the first map is injective exactly when the \(u_\rho\) span
\(N_\RR\).
\end{proof}

The most striking part of the theorem is the surjectivity statement:

\[
\boxed{
\text{Every Weil divisor class has a torus-invariant representative.}
}
\]

A general divisor need not itself be invariant. It is only its class
that can always be represented by a sum of boundary divisors.

\begin{pauseandthink}
Suppose \(\Sigma=\{\{0\}\}\) in a nonzero lattice \(N\). Then
\[
X_\Sigma=T_N.
\]
What are
\[
\Div_T(X_\Sigma)
\qquad\text{and}\qquad
\Cl(X_\Sigma)?
\]
Why is the map
\[
M\longrightarrow\Div_T(X_\Sigma)
\]
not injective?
\end{pauseandthink}

\subsection{A matrix for the class group}

Choose a basis
\[
m_1,\ldots,m_n
\]
of \(M\), and order the rays
\[
\Sigma(1)=\{\rho_1,\ldots,\rho_r\}.
\]
Write the primitive generators as column vectors
\[
u_i=u_{\rho_i}\in\ZZ^n.
\]

Under the identifications
\[
M\simeq\ZZ^n,
\qquad
\Div_T(X_\Sigma)\simeq\ZZ^r,
\]
the divisor map is represented by the \(r\times n\) matrix
\[
A=
\begin{pmatrix}
---u_1^{\,T}---\\
---u_2^{\,T}---\\
\vdots\\
---u_r^{\,T}---
\end{pmatrix}.
\]
Indeed, its \(i\)-th entry is
\[
\langle m,u_i\rangle.
\]
Consequently,
\[
\Cl(X_\Sigma)
\simeq
\coker(A).
\]

\begin{observation}
The Smith normal form of \(A\) separates the free and torsion parts of
the class group. When the rays span \(N_\RR\),
\[
\rank\Cl(X_\Sigma)
=
|\Sigma(1)|-\rank N.
\]
Torsion can still occur: it measures a finite-index lattice
phenomenon rather than a failure of the rays to span the real vector
space.
\end{observation}

For hand calculations, it is often faster to write down the
relations
\[
\sum_{\rho}
\langle m_i,u_\rho\rangle[D_\rho]
=
0
\]
coming from a basis \(m_1,\ldots,m_n\) of \(M\).

\subsection{Three familiar complete surfaces}

\begin{example}[The projective plane]
For \(\PP^2\), take
\[
u_0=-e_1-e_2,
\qquad
u_1=e_1,
\qquad
u_2=e_2.
\]
The characters \(m_1,m_2\) give the relations
\[
-D_0+D_1\sim0,
\qquad
-D_0+D_2\sim0.
\]
Thus
\[
[D_0]=[D_1]=[D_2].
\]
Writing their common class as \(H\), we obtain
\[
\Cl(\PP^2)
\simeq
\ZZ H.
\]
The three coordinate lines are different divisors but represent the
same divisor class.
\end{example}

\begin{example}[The product \(\PP^1\times\PP^1\)]
Order the rays
\[
e_1,\quad e_2,\quad -e_1,\quad -e_2,
\]
and denote the corresponding divisors by
\[
D_1,D_2,D_3,D_4.
\]
The two character relations are
\[
D_1-D_3\sim0,
\qquad
D_2-D_4\sim0.
\]
Hence
\[
\Cl(\PP^1\times\PP^1)
\simeq
\ZZ[D_1]\oplus\ZZ[D_2].
\]

The two generators have a geometric interpretation:
\[
[D_1]=[D_3]
\]
is the class of a fiber of one projection, while
\[
[D_2]=[D_4]
\]
is the class of a fiber of the other.
\end{example}

These two examples illustrate the rank count:
\[
3-2=1
\quad\text{for }\PP^2,
\qquad
4-2=2
\quad\text{for }\PP^1\times\PP^1.
\]

\subsection{The blowup of \texorpdfstring{\(\PP^2\)}{P2}}

Blow up the fixed point corresponding to
\[
\Cone{e_1,e_2}.
\]
The resulting fan has primitive ray generators
\[
u_0=-e_1-e_2,
\qquad
u_1=e_1,
\qquad
u_E=e_1+e_2,
\qquad
u_2=e_2.
\]
Let the corresponding divisors be
\[
D_0,\qquad D_1,\qquad E,\qquad D_2.
\]

\begin{figure}[ht]
\centering
\begin{tikzpicture}[scale=1.05]
  \fill[DictionaryBlue!9]
    (0,0)--(2.5,0)--(1.25,1.25)--cycle;
  \fill[PredictionTeal!10]
    (0,0)--(1.25,1.25)--(0,2.5)--cycle;
  \fill[SurfaceGreen!9]
    (0,0)--(0,2.5)--(-1.8,-1.8)--cycle;
  \fill[ExampleGold!10]
    (0,0)--(-1.8,-1.8)--(2.5,0)--cycle;

  \draw[very thick,DictionaryBlue,-{Latex[length=2.3mm]}]
    (0,0)--(2.85,0);
  \draw[very thick,PredictionTeal,-{Latex[length=2.3mm]}]
    (0,0)--(2.2,2.2);
  \draw[very thick,SurfaceGreen,-{Latex[length=2.3mm]}]
    (0,0)--(0,2.85);
  \draw[very thick,ExampleGold,-{Latex[length=2.3mm]}]
    (0,0)--(-2.1,-2.1);

  \node[DictionaryBlue,below] at (2.65,0)
    {\(D_1\sim H-E\)};
  \node[PredictionTeal,above right] at (1.92,1.92)
    {\(E\)};
  \node[SurfaceGreen,left] at (0,2.65)
    {\(D_2\sim H-E\)};
  \node[ExampleGold,above left] at (-1.9,-1.9)
    {\(D_0\sim H\)};

  \node[align=center] at (0.5,3.6)
    {one ray adds one generator\\
     to the divisor class group};
\end{tikzpicture}
\caption{Divisor classes on \(\Bl_p(\PP^2)\). The new ray contributes
the exceptional class \(E\); the two boundary lines through the
blown-up point have class \(H-E\).}
\label{fig:blowup-p2-divisor-classes}
\end{figure}

The character relations are
\[
-D_0+D_1+E\sim0,
\qquad
-D_0+E+D_2\sim0.
\]
Therefore
\[
[D_1]=[D_0]-[E],
\qquad
[D_2]=[D_0]-[E].
\]
If
\[
H=[D_0],
\]
then
\[
\Cl(\Blp{\PP^2}{p})
\simeq
\ZZ H\oplus\ZZ E.
\]

Geometrically, \(D_0\) is the strict transform of the coordinate line
that misses the blown-up point, so its class is the pullback of a line.
The divisors \(D_1,D_2\) are strict transforms of lines through the
point, hence their class is
\[
H-E.
\]

\begin{surfacelaboratory}
Use the fan and the principal-divisor formula to recover the class
group of \(\Blp{\PP^2}{p}\).

\begin{enumerate}[label=(\alph*),wide=0pt]
\item Write the \(4\times2\) matrix representing
\[
M\longrightarrow\Div_T(X).
\]

\item Verify directly that its cokernel is free of rank two.

\item Explain why adding one ray to the fan of \(\PP^2\) increases the
rank of the class group by one.

\item Compare the relation
\[
[D_1]=H-E
\]
with the total-transform formula
\[
\pi^*L=\widetilde L+E
\]
for a line \(L\) through the blown-up point.
\end{enumerate}
\end{surfacelaboratory}

\subsection{A singular comparison: torsion}

Let
\[
X_n
=
\AAA^2/\tfrac1n(1,n-1)
\]
be the \(A_{n-1}\) cyclic quotient from the Chapter~3 challenge. Its
fan consists of the first-quadrant cone in
\[
N
=
\ZZ^2+
\ZZ\left(\frac1n(1,n-1)\right).
\]
The two primitive ray generators are still
\[
e_1,e_2,
\]
but they generate a sublattice
\[
\ZZ e_1+\ZZ e_2\subseteq N
\]
of index \(n\).

The class-group matrix has full real rank, but its image is not
saturated. Its Smith normal form gives
\[
\Cl(X_n)
\simeq
\ZZ/n\ZZ.
\]

Thus the class group detects the finite quotient hidden in the
lattice. The contrast is worth remembering:

\[
\begin{array}{c|c}
\text{extra rays}
&
\text{free divisor classes}
\\
\hline
\text{finite-index lattice change}
&
\text{torsion divisor classes}.
\end{array}
\]

\begin{pauseandthink}
Why does the resolution fan of \(X_n\), which inserts \(n-1\) new
rays and is smooth, have a torsion-free class group?

Use only the facts that adjacent primitive generators form lattice
bases and that the full collection of ray generators therefore spans
\(N\) as an abelian group.
\end{pauseandthink}

The class group treats Weil divisors modulo rational equations. The
next question is more local: when can a divisor itself be described
by one equation on each affine chart?

\begin{guidingquestion}
For an invariant divisor
\[
D=\sum_\rho a_\rho D_\rho,
\]
when can we choose, on every \(U_\sigma\), a character whose divisor
agrees with \(D\) along all rays of \(\sigma\)?
\end{guidingquestion}


\section{Cartier divisors and local equations}
\label{sec:cartier-divisors-local-equations}
\index{Cartier divisor}
\index{Cartier data}
\index{local equation!of a divisor}
\index{Picard group}

A Weil divisor records integer multiplicities along codimension-one
subvarieties. A Cartier divisor asks for more: those multiplicities
must arise from compatible local equations.

\subsection{From multiplicities to local equations}

\begin{agcheatsheet}
Let \(X\) be an irreducible variety.

A \emph{Cartier divisor} is described by an open cover
\[
X=\bigcup_i U_i
\]
and nonzero rational functions
\[
f_i\in\Bbbk(X)^\times
\]
such that
\[
f_i/f_j
\]
is regular and nowhere zero on every overlap \(U_i\cap U_j\).

The function \(f_i\) is a \emph{local equation} for the divisor on
\(U_i\). On a normal variety, the local equations determine a Weil
divisor by
\[
D|_{U_i}=\operatorname{div}(f_i)|_{U_i}.
\]

Thus a Cartier divisor is a Weil divisor that is locally principal.
A global principal divisor is the special case in which one rational
function works on all of \(X\).
\end{agcheatsheet}

We will identify a Cartier divisor with its associated Weil divisor.
In particular, an invariant Weil divisor
\[
D=\sum_{\rho\in\Sigma(1)}a_\rho D_\rho
\]
is Cartier precisely when it can be cut out locally by one rational
function on each affine chart.

The toric setting is especially rigid: invariant local equations can
be chosen to be characters.

\begin{proposition}[Invariant Cartier divisors on an affine chart]
\label{prop:affine-invariant-cartier-principal}
Let \(\sigma\subseteq N_\RR\) be a strongly convex rational
polyhedral cone. Every torus-invariant Cartier divisor on
\[
U_\sigma
\]
is the divisor of a character. Consequently,
\[
\Pic(U_\sigma)=0.
\]
\end{proposition}

\noindent
\textit{Reference.}
This is Cox--Little--Schenck, Proposition~4.2.2.

\begin{proof}[Conceptual proof sketch]
The affine chart \(U_\sigma\) contains a unique closed torus orbit.
Begin with a local equation near that orbit. Because the divisor is
torus-invariant, translating the equation by the torus changes it only
by a unit.

The key step is to replace the local equation by one of its torus
eigencomponents. Eigenfunctions for the torus action are characters,
so one obtains a character having the same divisor near the closed
orbit. Every invariant boundary divisor of \(U_\sigma\) meets that
orbit, forcing equality of the two divisors on the entire affine
chart.

For a general, not necessarily effective, Cartier divisor, add a
sufficiently positive principal character divisor, apply the effective
case, and subtract it again.
\end{proof}

The technical point in this proposition is not that every Cartier
divisor on an affine variety is principal—that statement is false in
general. It is the combination of affineness and torus invariance that
forces a character equation.

\subsection{Cartier data on a fan}

Let
\[
D=\sum_{\rho\in\Sigma(1)}a_\rho D_\rho.
\]
On the affine chart \(U_\sigma\), only the boundary divisors
corresponding to rays of \(\sigma\) appear:
\[
D|_{U_\sigma}
=
\sum_{\rho\in\sigma(1)}a_\rho D_\rho.
\]

If this restriction is the divisor of \(\character{-m_\sigma}\), then
the principal-divisor formula gives
\[
a_\rho
=
\langle -m_\sigma,u_\rho\rangle
=
-\langle m_\sigma,u_\rho\rangle.
\]
This explains the signs in the following criterion.

\begin{theorem}[Toric Cartier criterion]
\label{thm:toric-cartier-criterion}
Let
\[
D=\sum_{\rho\in\Sigma(1)}a_\rho D_\rho
\]
be a torus-invariant Weil divisor. The following are equivalent:

\begin{enumerate}[label=(\alph*),wide=0pt]
\item \(D\) is Cartier.

\item \(D|_{U_\sigma}\) is principal for every \(\sigma\in\Sigma\).

\item For every \(\sigma\in\Sigma\), there exists
\[
m_\sigma\in M
\]
such that
\[
\langle m_\sigma,u_\rho\rangle=-a_\rho
\qquad
\text{for every }\rho\in\sigma(1).
\]

\item The condition in \textup{(c)} holds for every maximal cone of
\(\Sigma\).
\end{enumerate}

When these conditions hold, the local equation on \(U_\sigma\) is
\[
\character{-m_\sigma}.
\]
Moreover:
\[
m_\sigma
\text{ is unique modulo }
M(\sigma):=\sigma^\perp\cap M,
\]
and whenever \(\tau\preceq\sigma\),
\[
m_\sigma\equiv m_\tau\pmod{M(\tau)}.
\]
The collection
\[
\{m_\sigma\}_{\sigma\in\Sigma}
\]
is called the \emph{Cartier data} of \(D\).
\end{theorem}

\noindent
\textit{Reference.}
This is Cox--Little--Schenck, Theorem~4.2.8.

\begin{proof}
A Cartier divisor restricts to an invariant Cartier divisor on every
affine chart. By
\cref{prop:affine-invariant-cartier-principal}, that restriction is
the divisor of a character. The principal-divisor formula shows that
the character can be written as \(\character{-m_\sigma}\) precisely
when
\[
\langle m_\sigma,u_\rho\rangle=-a_\rho
\]
for all rays of \(\sigma\). This proves the equivalence of
\textup{(a)}, \textup{(b)}, and \textup{(c)}.

It is enough to work on maximal cones because their affine charts
cover \(X_\Sigma\), and any character equation for a maximal cone
restricts to every face.

If \(m_\sigma\) and \(m'_\sigma\) give the same coefficients, then
\[
\langle m_\sigma-m'_\sigma,u_\rho\rangle=0
\]
on every ray of \(\sigma\), hence on all of \(\sigma\). Thus
\[
m_\sigma-m'_\sigma\in\sigma^\perp\cap M.
\]
The compatibility along faces follows by applying this uniqueness
statement to the restricted divisor on \(U_\tau\).
\end{proof}

\begin{warningbox}
The coefficient convention is
\[
D=\sum_\rho a_\rho D_\rho,
\qquad
\langle m_\sigma,u_\rho\rangle=-a_\rho,
\]
because the local equation is
\[
\character{-m_\sigma}.
\]
It is easy to lose this minus sign when passing between a divisor and
its Cartier data.
\end{warningbox}

If \(\sigma\) is full-dimensional, then
\[
M(\sigma)=0,
\]
so \(m_\sigma\) is uniquely determined. For a lower-dimensional cone,
one may add any character that is trivial on the subtorus directions
seen by \(\sigma\).

\begin{pauseandthink}
Suppose \(\tau\) is a ray. Why can two different elements of \(M\)
give the same local equation data along \(D_\tau\)?

Describe the ambiguity geometrically using
\[
M(\tau)=\tau^\perp\cap M.
\]
\end{pauseandthink}

\subsection{A coordinate line on \texorpdfstring{\(\PP^2\)}{P2}}

Let the primitive ray generators be
\[
u_0=-e_1-e_2,
\qquad
u_1=e_1,
\qquad
u_2=e_2,
\]
and let
\[
D_0=\{z_0=0\}.
\]
The maximal cones are
\[
\sigma_0=\Cone{u_1,u_2},
\qquad
\sigma_1=\Cone{u_0,u_2},
\qquad
\sigma_2=\Cone{u_0,u_1}.
\]

For
\[
D_0=1D_0+0D_1+0D_2,
\]
the Cartier data and local equations are:

\begin{center}
\small
\renewcommand{\arraystretch}{1.25}
\begin{tabularx}{0.92\linewidth}{@{}
  >{\raggedright\arraybackslash}p{0.18\linewidth}
  >{\raggedright\arraybackslash}p{0.25\linewidth}
  >{\raggedright\arraybackslash}p{0.22\linewidth}
  >{\raggedright\arraybackslash}X@{}}
\toprule
\textbf{Chart}
&
\textbf{Rays in the cone}
&
\textbf{Cartier datum}
&
\textbf{Local equation}
\\
\midrule
\(U_{\sigma_0}\)
&
\(u_1,u_2\)
&
\(m_{\sigma_0}=0\)
&
\(1\)
\\
\(U_{\sigma_1}\)
&
\(u_0,u_2\)
&
\(m_{\sigma_1}=m_1\)
&
\(\character{-m_1}=z_0/z_1\)
\\
\(U_{\sigma_2}\)
&
\(u_0,u_1\)
&
\(m_{\sigma_2}=m_2\)
&
\(\character{-m_2}=z_0/z_2\)
\\
\bottomrule
\end{tabularx}
\end{center}

On \(U_{\sigma_0}\), the divisor \(D_0\) is absent, so the local
equation is a unit. On the other two charts, the displayed ratios are
the familiar equations of the coordinate line.

On an overlap, the quotient of two local equations is a character
whose pairing with every ray of the common face is zero. It is
therefore a unit on that overlap.

\begin{surfacelaboratory}
Repeat the calculation for
\[
D_1=\{z_1=0\}.
\]

\begin{enumerate}[label=(\alph*),wide=0pt]
\item Find the three elements
\[
m_{\sigma_0},m_{\sigma_1},m_{\sigma_2}.
\]

\item Write the local equations in homogeneous coordinates.

\item Verify explicitly that the ratios of the equations are units on
the pairwise overlaps.
\end{enumerate}
\end{surfacelaboratory}

\subsection{A Weil divisor that is not Cartier}

Return to the cyclic quotient
\[
X_n
=
\AAA^2/\tfrac1n(1,n-1).
\]
Its cone is the first quadrant, but its character lattice is
\[
M
=
\left\{
(a,b)\in\ZZ^2:
a-b\equiv0\pmod n
\right\}.
\]
Let \(D_1,D_2\) be the invariant prime divisors corresponding to the
two coordinate rays.

Since \(X_n\) is affine, an invariant Cartier divisor must be the
divisor of a character. To have
\[
\operatorname{div}(\character{m})=D_1,
\]
we would need
\[
m=(1,0).
\]
But
\[
(1,0)\notin M
\qquad
(n>1),
\]
so \(D_1\) is not Cartier.

On the other hand,
\[
(n,0)\in M,
\]
and
\[
\operatorname{div}(\character{(n,0)})
=
nD_1.
\]
Thus \(D_1\) is not Cartier, but a positive multiple of it is.

\begin{definition}
A Weil divisor \(D\) is \emph{\(\QQ\)-Cartier} if some positive
integer multiple of \(D\) is Cartier.
\end{definition}

In this example, \(D_1\) is \(\QQ\)-Cartier of Cartier index \(n\).
The same is true of \(D_2\).

This gives a sharp distinction:
\[
\Pic(X_n)=0,
\qquad
\Cl(X_n)\simeq\ZZ/n\ZZ.
\]
The class group sees Weil divisors that do not admit local equations;
the Picard group sees only Cartier classes.

\subsection{The exceptional divisor on the blowup}

Let
\[
X=\Blp{\PP^2}{p}
\]
with ray generators
\[
u_0=-e_1-e_2,\qquad
u_1=e_1,\qquad
u_E=e_1+e_2,\qquad
u_2=e_2.
\]
We compute local equations for the exceptional divisor \(E=D_{\rho_E}\).

\begin{figure}[ht]
\centering
\begin{tikzpicture}[scale=1.02]
  \fill[DictionaryBlue!8]
    (0,0)--(2.55,0)--(1.28,1.28)--cycle;
  \fill[PredictionTeal!10]
    (0,0)--(1.28,1.28)--(0,2.55)--cycle;
  \fill[SurfaceGreen!8]
    (0,0)--(0,2.55)--(-1.85,-1.85)--cycle;
  \fill[ExampleGold!9]
    (0,0)--(-1.85,-1.85)--(2.55,0)--cycle;

  \draw[very thick,DictionaryBlue,-{Latex[length=2.3mm]}]
    (0,0)--(2.9,0);
  \draw[very thick,PredictionTeal,-{Latex[length=2.3mm]}]
    (0,0)--(2.25,2.25);
  \draw[very thick,SurfaceGreen,-{Latex[length=2.3mm]}]
    (0,0)--(0,2.9);
  \draw[very thick,ExampleGold,-{Latex[length=2.3mm]}]
    (0,0)--(-2.15,-2.15);

  \node[DictionaryBlue,below] at (2.68,0) {\(u_1\)};
  \node[PredictionTeal,above right] at (2.0,2.0) {\(u_E\)};
  \node[SurfaceGreen,left] at (0,2.68) {\(u_2\)};
  \node[ExampleGold,above left] at (-1.95,-1.95) {\(u_0\)};

  \node[fill=DictionaryBlue!12,rounded corners,inner sep=2pt]
    at (1.75,0.48)
    {\(\character{-m_\sigma}=y\)};
  \node[fill=PredictionTeal!14,rounded corners,inner sep=2pt]
    at (0.48,1.75)
    {\(\character{-m_\sigma}=x\)};
  \node[fill=SurfaceGreen!12,rounded corners,inner sep=2pt]
    at (-0.55,0.85)
    {\(\character{-m_\sigma}=1\)};
  \node[fill=ExampleGold!14,rounded corners,inner sep=2pt]
    at (0.5,-0.78)
    {\(\character{-m_\sigma}=1\)};

  \node[align=center] at (0.5,3.65)
    {local equations for the exceptional divisor \(E\)};
\end{tikzpicture}
\caption{Cartier data for the exceptional divisor. On the two charts
meeting \(E\), its local equations are \(y\) and \(x\); on the other
two charts, \(E\) is absent and the equation is \(1\).}
\label{fig:exceptional-divisor-cartier-data}
\end{figure}

For the four maximal cones, we may take
\[
\begin{array}{c|c|c}
\text{maximal cone}
&
m_\sigma
&
\character{-m_\sigma}
\\
\hline
\Cone{u_0,u_1}
&
0
&
1
\\
\Cone{u_1,u_E}
&
-m_2
&
y
\\
\Cone{u_E,u_2}
&
-m_1
&
x
\\
\Cone{u_2,u_0}
&
0
&
1.
\end{array}
\]
Indeed, on the cone \(\Cone{u_1,u_E}\),
\[
\langle -m_2,u_1\rangle=0,
\qquad
\langle -m_2,u_E\rangle=-1,
\]
and similarly on \(\Cone{u_E,u_2}\).

The two nontrivial equations agree up to a unit on their overlap:
\[
\frac{x}{y}
\]
is invertible there. This is the local-equation version of the gluing
that produced the blowup.

\begin{observation}
The calculation succeeds because the primitive generators of each
maximal cone form a lattice basis. Given arbitrary integer
coefficients on the two rays, there is a unique integral linear
function taking the required values.

This is the local reason smooth toric varieties have no distinction
between Weil and Cartier divisors. We will return to the precise
criterion later in this chapter.
\end{observation}

\subsection{Cartier classes}

Cartier divisors modulo principal divisors form the Picard group:
\[
\Pic(X)
=
\operatorname{CDiv}(X)/
\Prin(X).
\]
Equivalently, this is the group of line bundles on \(X\); that
interpretation will become more prominent in the next chapter.

Every Cartier class on a toric variety has an invariant
representative. Consequently, there is an exact sequence
\[
M
\longrightarrow
\operatorname{CDiv}_T(X_\Sigma)
\longrightarrow
\Pic(X_\Sigma)
\longrightarrow
0.
\]

\noindent
\textit{Reference.}
This is Cox--Little--Schenck, Theorem~4.2.1.

The difference between
\[
\Pic(X_\Sigma)
\qquad\text{and}\qquad
\Cl(X_\Sigma)
\]
measures the failure of Weil divisors to admit compatible local
equations.

The Cartier data are a collection of integral linear functions, one
on each cone, compatible on common faces. There should be a more
economical way to package them.

\begin{guidingquestion}
Can the collection
\[
\{m_\sigma\}_{\sigma\in\Sigma}
\]
be assembled into one function on \(|\Sigma|\) that is linear on each
cone and remembers the coefficients of the divisor on every ray?
\end{guidingquestion}


\section{Support functions}
\label{sec:support-functions}
\index{support function}
\index{piecewise-linear function}
\index{linear equivalence!global linear function}

The Cartier data
\[
\{m_\sigma\}_{\sigma\in\Sigma}
\]
assign one integral linear function to every cone. Compatibility says
that these linear functions agree wherever the cones meet.

This collection can therefore be viewed as one function on the entire
support of the fan.

\subsection{One function, linear on every cone}

\begin{definition}
\label{def:support-function}
Let \(\Sigma\) be a fan in \(N_\RR\).

A \emph{support function} is a function
\[
\psi:|\Sigma|\longrightarrow\RR
\]
whose restriction to every cone \(\sigma\in\Sigma\) is linear.

The support function is \emph{integral} if
\[
\psi(|\Sigma|\cap N)\subseteq\ZZ.
\]
We write
\[
\operatorname{SF}(\Sigma)
\]
for the group of support functions and
\[
\operatorname{SF}(\Sigma,N)
\]
for the subgroup of integral support functions.
\end{definition}

Because adjacent linear pieces agree on their common face, a support
function is automatically continuous. It is also homogeneous:
\[
\psi(tu)=t\psi(u)
\qquad
(t\geq0).
\]
Thus ``piecewise linear'' here means linear, rather than affine
linear, on each cone.

For a rational cone \(\sigma\), integrality is equivalent to the
existence of some \(m_\sigma\in M\) such that
\[
\psi(u)=\langle m_\sigma,u\rangle
\qquad
(u\in\sigma).
\]
The element \(m_\sigma\) is determined only modulo
\[
M(\sigma)=\sigma^\perp\cap M,
\]
exactly as in the Cartier criterion.

\noindent
\textit{Reference.}
This is Cox--Little--Schenck, Definition~4.2.11.

\subsection{Cartier divisors and integral support functions}

Let
\[
D=\sum_{\rho\in\Sigma(1)}a_\rho D_\rho
\]
be a torus-invariant Cartier divisor with Cartier data
\[
\{m_\sigma\}_{\sigma\in\Sigma}.
\]
Define
\[
\psi_D(u)
=
\langle m_\sigma,u\rangle
\qquad
\text{when }u\in\sigma.
\]

\begin{theorem}[The support-function dictionary]
\label{thm:support-function-dictionary}
The function \(\psi_D\) is a well-defined integral support function.
Its values on the primitive ray generators recover the divisor:
\[
\psi_D(u_\rho)=-a_\rho,
\]
and hence
\[
\boxed{
D
=
-\sum_{\rho\in\Sigma(1)}
\psi_D(u_\rho)D_\rho.
}
\]

Moreover, the assignment
\[
D\longmapsto\psi_D
\]
is an isomorphism of abelian groups
\[
\operatorname{CDiv}_T(X_\Sigma)
\simeq
\operatorname{SF}(\Sigma,N).
\]
\end{theorem}

\noindent
\textit{Reference.}
This is Cox--Little--Schenck, Theorem~4.2.12.

\begin{proof}
If \(u\in\sigma\cap\sigma'\), compatibility of the Cartier data gives
\[
m_\sigma-m_{\sigma'}
\in
(\sigma\cap\sigma')^\perp\cap M.
\]
Therefore
\[
\langle m_\sigma,u\rangle
=
\langle m_{\sigma'},u\rangle,
\]
so \(\psi_D\) is well defined. It is integral and linear on each cone
because every \(m_\sigma\) lies in \(M\).

For a ray \(\rho\preceq\sigma\), the Cartier criterion gives
\[
\psi_D(u_\rho)
=
\langle m_\sigma,u_\rho\rangle
=
-a_\rho.
\]
Thus the ray values recover \(D\), proving injectivity.

Conversely, let
\[
\psi\in\operatorname{SF}(\Sigma,N).
\]
On each cone, choose \(m_\sigma\in M\) with
\[
\psi|_\sigma=\langle m_\sigma,-\rangle.
\]
The pieces agree on common faces, so these \(m_\sigma\) are compatible
Cartier data. The divisor
\[
D_\psi
=
-\sum_{\rho\in\Sigma(1)}
\psi(u_\rho)D_\rho
\]
has support function \(\psi\). This proves surjectivity.
\end{proof}

The theorem compresses the Cartier criterion into a picture:
\[
\begin{array}{c@{\qquad\longleftrightarrow\qquad}c}
\text{local characters }m_\sigma
&
\text{linear pieces of }\psi_D
\\[1mm]
\langle m_\sigma,u_\rho\rangle=-a_\rho
&
\psi_D(u_\rho)=-a_\rho
\\[1mm]
\text{compatibility on overlaps}
&
\text{one well-defined function on }|\Sigma|.
\end{array}
\]

\begin{warningbox}
An arbitrary assignment of integers to the rays need not extend to a
support function. On every cone, the assigned values must come from
one linear function.

This is exactly the obstruction to the corresponding invariant Weil
divisor being Cartier.
\end{warningbox}

\subsection{The support function of a line in \texorpdfstring{\(\PP^2\)}{P2}}

Consider
\[
D_0=\{z_0=0\}\subseteq\PP^2.
\]
With
\[
u_0=-e_1-e_2,\qquad
u_1=e_1,\qquad
u_2=e_2,
\]
the ray values are
\[
\psi_{D_0}(u_0)=-1,
\qquad
\psi_{D_0}(u_1)=0,
\qquad
\psi_{D_0}(u_2)=0.
\]

The Cartier data from the previous section give
\[
\psi_{D_0}(x,y)
=
\begin{cases}
0,
&
(x,y)\in\Cone{e_1,e_2},
\\[1mm]
x,
&
(x,y)\in\Cone{-e_1-e_2,e_2},
\\[1mm]
y,
&
(x,y)\in\Cone{e_1,-e_1-e_2}.
\end{cases}
\]

\begin{figure}[ht]
\centering
\begin{tikzpicture}[scale=0.88]
  % Fan with values and formulas
  \begin{scope}[xshift=-5.15cm,scale=0.82]
    \fill[DictionaryBlue!9]
      (0,0)--(2.5,0)--(0,2.5)--cycle;
    \fill[SurfaceGreen!9]
      (0,0)--(0,2.5)--(-1.85,-1.85)--cycle;
    \fill[ExampleGold!11]
      (0,0)--(-1.85,-1.85)--(2.5,0)--cycle;

    \draw[very thick,DictionaryBlue,-{Latex[length=2.3mm]}]
      (0,0)--(2.85,0);
    \draw[very thick,SurfaceGreen,-{Latex[length=2.3mm]}]
      (0,0)--(0,2.85);
    \draw[very thick,ExampleGold,-{Latex[length=2.3mm]}]
      (0,0)--(-2.15,-2.15);

    \node[DictionaryBlue,below] at (2.62,0)
      {\(\psi(u_1)=0\)};
    \node[SurfaceGreen,left] at (0,2.62)
      {\(\psi(u_2)=0\)};
    \node[ExampleGold,above left] at (-1.95,-1.95)
      {\(\psi(u_0)=-1\)};

    \node[fill=DictionaryBlue!13,rounded corners,inner sep=2pt]
      at (0.95,0.95) {\(\psi=0\)};
    \node[fill=SurfaceGreen!13,rounded corners,inner sep=2pt]
      at (-0.55,0.8) {\(\psi=x\)};
    \node[fill=ExampleGold!15,rounded corners,inner sep=2pt]
      at (0.75,-0.55) {\(\psi=y\)};

    \node[font=\bfseries] at (0.3,3.55)
      {values and linear pieces};
  \end{scope}

  % A clearer 3D graph of the PL function
  \begin{scope}[
    xshift=3.55cm,
    scale=1.13,
    x={(0.92cm,-0.28cm)},
    y={(0.42cm,0.32cm)},
    z={(0cm,0.86cm)}
  ]
    \coordinate (O) at (0,0,0);
    \coordinate (A) at (2.7,0,0);
    \coordinate (B) at (0,2.7,0);
    \coordinate (Cbase) at (-1.8,-1.8,0);
    \coordinate (C) at (-1.8,-1.8,-1.8);

    % Coordinate axes and projected base fan
    \draw[->,black!60] (O)--(3.25,0,0)
      node[below right] {\(x\)};
    \draw[->,black!60] (O)--(0,3.45,0)
      node[above] {\(y\)};
    \draw[->,black!60] (O)--(0,0,2.35)
      node[right] {\(\psi\)};

    \draw[black!35,dashed] (O)--(Cbase);
    \draw[black!35,dotted] (Cbase)--(C);
    \node[black!55,below left] at (Cbase)
      {\(u_0\) in the base};

    % Three planar facets of the graph
    \filldraw[
      fill=DictionaryBlue!18,
      fill opacity=0.82,
      draw=DictionaryBlue!85,
      very thick
    ] (O)--(A)--(B)--cycle;

    \filldraw[
      fill=SurfaceGreen!18,
      fill opacity=0.82,
      draw=SurfaceGreen!85,
      very thick
    ] (O)--(B)--(C)--cycle;

    \filldraw[
      fill=ExampleGold!22,
      fill opacity=0.82,
      draw=ExampleGold!90,
      very thick
    ] (O)--(C)--(A)--cycle;

    % Grid lines make each facet visibly planar
    \draw[DictionaryBlue!65]
      (1.35,0,0)--(0,1.35,0);
    \draw[DictionaryBlue!45]
      (0.68,0,0)--(0,0.68,0);

    \draw[SurfaceGreen!65]
      (0,1.35,0)--(-0.9,-0.9,-0.9);
    \draw[SurfaceGreen!45]
      (0,0.68,0)--(-0.45,-0.45,-0.45);

    \draw[ExampleGold!70]
      (-0.9,-0.9,-0.9)--(1.35,0,0);
    \draw[ExampleGold!50]
      (-0.45,-0.45,-0.45)--(0.68,0,0);

    % Lifted ray directions and labels
    \fill (O) circle (1.7pt);
    \fill (A) circle (1.7pt);
    \fill (B) circle (1.7pt);
    \fill (C) circle (1.7pt);

    \node[below left] at (C) {\((u_0,-1)\)};

    \node[
      fill=DictionaryBlue!12,
      rounded corners,
      inner sep=2pt
    ] at (0.9,0.9,0.08) {\(\psi=0\)};

    \node[
      fill=SurfaceGreen!12,
      rounded corners,
      inner sep=2pt
    ] at (-0.62,0.3,-0.57) {\(\psi=x\)};

    \node[
      fill=ExampleGold!16,
      rounded corners,
      inner sep=2pt
    ] at (0.35,-0.62,-0.57) {\(\psi=y\)};

      \end{scope}
\end{tikzpicture}
\caption{The support function of the coordinate line \(D_0\). The
right-hand panel shows a truncated portion of its graph in
\(N_\RR\times\RR\). Its three facets lie in the planes
\(\psi=0\), \(\psi=x\), and \(\psi=y\).}
\label{fig:p2-support-function-line}
\end{figure}

The right-hand picture is the graph
\[
\Gamma_{\psi_{D_0}}
=
\{(u,\psi_{D_0}(u)):u\in N_\RR\}
\subseteq
N_\RR\times\RR.
\]
Only a bounded triangular portion of each unbounded planar piece is
drawn. Over \(\Cone{u_1,u_2}\) the graph is horizontal; over
\(\Cone{u_0,u_2}\) it lies in the plane \(\psi=x\); and over
\(\Cone{u_1,u_0}\) it lies in the plane \(\psi=y\). The three planes
meet along the lifted common rays.

The bending of the graph records that \(D_0\) is not principal. A
principal divisor would produce one global linear function, with no
change of slope from cone to cone.

\begin{pauseandthink}
Use the ray values in
\cref{fig:p2-support-function-line} to recover the divisor \(D_0\)
from
\[
D=-\sum_\rho\psi_D(u_\rho)D_\rho.
\]

Then verify directly that the three displayed formulas agree on every
common ray.
\end{pauseandthink}

\subsection{Linear equivalence means changing the global tilt}

Let \(D\) be an invariant Cartier divisor and let \(m\in M\). Since
\[
\operatorname{div}(\character{m})
=
\sum_\rho\langle m,u_\rho\rangle D_\rho,
\]
we obtain
\[
\psi_{D+\operatorname{div}(\character{m})}(u)
=
\psi_D(u)-\langle m,u\rangle.
\]

Thus adding a principal divisor changes the support function by one
global integral linear function. It does not alter how the linear
pieces bend relative to one another.

Consequently, the Picard group has the presentation
\[
M
\longrightarrow
\operatorname{SF}(\Sigma,N)
\longrightarrow
\Pic(X_\Sigma)
\longrightarrow
0,
\]
where
\[
m\longmapsto
\bigl(u\longmapsto-\langle m,u\rangle\bigr).
\]
Equivalently,
\[
\Pic(X_\Sigma)
\simeq
\frac{\operatorname{SF}(\Sigma,N)}
{\{\text{restrictions of global integral linear functions}\}}.
\]

\noindent
\textit{Reference.}
This is the support-function form of Cox--Little--Schenck,
Theorem~4.2.1, following Theorem~4.2.12.

\begin{example}[Why the coordinate lines of \(\PP^2\) are equivalent]
Let \(\psi_0\) and \(\psi_1\) be the support functions of \(D_0\) and
\(D_1\). Since
\[
D_1-D_0
=
\operatorname{div}(z_1/z_0)
=
\operatorname{div}(\character{m_1}),
\]
we have
\[
\psi_1(u)
=
\psi_0(u)-\langle m_1,u\rangle.
\]
Their graphs differ by a global tilt, so they represent the same
Picard class.
\end{example}

This gives a useful geometric interpretation:
\[
\boxed{
\text{a Cartier class remembers the bending of a support function,
not its overall tilt.}
}
\]

\subsection{Surface laboratory: the exceptional divisor}

Let
\[
X=\Blp{\PP^2}{p}
\]
and use the rays
\[
u_0=-e_1-e_2,\qquad
u_1=e_1,\qquad
u_E=e_1+e_2,\qquad
u_2=e_2.
\]
For the exceptional divisor \(E\), the support function satisfies
\[
\psi_E(u_E)=-1,
\qquad
\psi_E(u_0)=\psi_E(u_1)=\psi_E(u_2)=0.
\]

\begin{surfacelaboratory}
\begin{enumerate}[label=(\alph*),wide=0pt]
\item Use the Cartier data from
\cref{fig:exceptional-divisor-cartier-data} to show that
\[
\psi_E(x,y)
=
\begin{cases}
0,
&
\Cone{u_0,u_1}\text{ and }\Cone{u_2,u_0},
\\
-y,
&
\Cone{u_1,u_E},
\\
-x,
&
\Cone{u_E,u_2}.
\end{cases}
\]

\item Check that the four pieces agree on their common rays.

\item Sketch the graph of \(\psi_E\). Where does it bend?

\item Explain why the newly inserted ray is exactly where the support
function acquires a new linear piece.

\item Compute the support function of the invariant divisor
\[
D_0-E
\]
by subtracting support functions. Which Picard class does this divisor
represent?
\end{enumerate}
\end{surfacelaboratory}

\begin{observation}
Support functions will return in the next chapter with additional
shape conditions. Convexity and strict convexity will translate into
global generation and ampleness, respectively, while the supporting
inequalities
\[
\langle m,u_\rho\rangle\geq-\!a_\rho
\]
will cut out a polyhedron in \(M_\RR\).

For now, the essential point is only the Cartier dictionary:
integral piecewise-linear functions encode compatible local
equations.
\end{observation}

A particularly important invariant divisor assigns the same
coefficient to every boundary component. Its Cartier condition asks
for all primitive ray generators of each cone to lie on a common
integral affine hyperplane.

\begin{guidingquestion}
What geometric divisor has coefficient \(-1\) on every
\(D_\rho\), and how does the existence of a support function with
\[
\psi(u_\rho)=1
\]
detect whether that divisor is Cartier?
\end{guidingquestion}


\section{The canonical divisor and singularity criteria}
\label{sec:canonical-divisor-singularity-criteria}
\index{canonical divisor!toric formula}
\index{factorial toric variety}
\index{Q-factorial toric variety}
\index{Gorenstein cone}
\index{Q-Gorenstein cone}
\index{height-one criterion}

The toric boundary is the reduced divisor
\[
\partial X_\Sigma
=
\sum_{\rho\in\Sigma(1)}D_\rho.
\]
The canonical divisor is its negative:
\[
K_{X_\Sigma}
=
-\partial X_\Sigma.
\]

This simple formula brings together every dictionary developed in the
chapter. It also lets us recognize several singularity properties
directly from the primitive ray generators.

\subsection{The canonical divisor is minus the boundary}

\begin{agcheatsheet}
Let \(X\) be a smooth irreducible variety of dimension \(n\). A
nonzero rational top differential form
\[
\omega
\]
has an order of vanishing along every prime divisor. Its divisor
\[
\operatorname{div}(\omega)
\]
is called a \emph{canonical divisor}.

Different rational top forms differ by a nonzero rational function,
so their divisors are linearly equivalent. Thus the canonical divisor
class
\[
[K_X]
\]
is intrinsic even though a particular representative is not.

For a normal variety, the smooth locus contains every codimension-one
point. One may therefore compute the canonical divisor on the smooth
locus and take the closures of its codimension-one components.
\end{agcheatsheet}

Choose a basis
\[
m_1,\ldots,m_n
\]
of \(M\), and write
\[
x_i=\character{m_i}.
\]
On the dense torus, consider the invariant logarithmic volume form
\[
\Omega_T
=
\frac{dx_1}{x_1}\wedge\cdots\wedge\frac{dx_n}{x_n}.
\]

\begin{theorem}[The canonical divisor of a toric variety]
\label{thm:canonical-divisor-toric}
For every normal toric variety \(X_\Sigma\),
\[
\boxed{
K_{X_\Sigma}
=
-\sum_{\rho\in\Sigma(1)}D_\rho
}
\]
is a torus-invariant canonical divisor. Equivalently,
\[
-K_{X_\Sigma}
=
\partial X_\Sigma.
\]
\end{theorem}

\noindent
\textit{Reference.}
This is Cox--Little--Schenck, Theorem~8.2.3.

\begin{proof}
The form \(\Omega_T\) is regular and nowhere zero on the dense torus,
so any zeros or poles occur along the boundary.

Fix a ray \(\rho\). Because its generator \(u_\rho\) is primitive, it
extends to a basis
\[
u_\rho=e_1,e_2,\ldots,e_n
\]
of \(N\). Let \(m_1,\ldots,m_n\) now denote the dual basis. On the ray
chart,
\[
U_\rho
\simeq
\AAA^1\times\Gm^{n-1},
\]
with
\[
D_\rho\cap U_\rho=V(x_1).
\]
Up to a nonzero constant, the invariant form becomes
\[
\Omega_T
=
\frac{dx_1}{x_1}\wedge
\frac{dx_2}{x_2}\wedge\cdots\wedge
\frac{dx_n}{x_n}.
\]
The factors \(x_2,\ldots,x_n\) are units along \(D_\rho\), while
\(dx_1/x_1\) has a simple pole. Therefore
\[
\ord_{D_\rho}(\Omega_T)=-1.
\]
This holds for every ray, proving the formula.
\end{proof}

\begin{observation}
The displayed divisor is one convenient canonical representative,
not necessarily the divisor of the most familiar volume form.

For example, on \(\AAA^2\),
\[
dx\wedge dy
\]
has divisor \(0\), while
\[
\frac{dx}{x}\wedge\frac{dy}{y}
=
\frac{1}{xy}\,dx\wedge dy
\]
has divisor
\[
-D_{\rho_1}-D_{\rho_2}.
\]
The two canonical divisors differ by
\[
\operatorname{div}(1/xy).
\]
\end{observation}

\subsection{Canonical classes of familiar surfaces}

The formula becomes especially transparent once the boundary classes
are known.

\begin{center}
\small
\renewcommand{\arraystretch}{1.3}
\begin{tabularx}{0.94\linewidth}{@{}
  >{\raggedright\arraybackslash}p{0.23\linewidth}
  >{\raggedright\arraybackslash}p{0.36\linewidth}
  >{\raggedright\arraybackslash}X@{}}
\toprule
\textbf{Surface}
&
\textbf{Boundary class}
&
\textbf{Canonical class}
\\
\midrule
\(\PP^2\)
&
\(D_0+D_1+D_2\sim3H\)
&
\(K_{\PP^2}\sim-3H\)
\\
\(\PP^1\times\PP^1\)
&
\(2F_1+2F_2\)
&
\(K\sim-2F_1-2F_2\)
\\
\(\Blp{\PP^2}{p}\)
&
\(D_0+D_1+E+D_2\sim3H-E\)
&
\(K\sim-3H+E\)
\\
\bottomrule
\end{tabularx}
\end{center}

For the blowup, recall
\[
[D_0]=H,
\qquad
[D_1]=[D_2]=H-E.
\]
Hence
\[
-K
=
D_0+D_1+E+D_2
\sim
H+(H-E)+E+(H-E)
=
3H-E.
\]

\begin{pauseandthink}
Compute the canonical class of the Hirzebruch surface
\(\mathbb F_r\) by summing its four invariant boundary divisors.

At this stage, it is enough to express the answer using any two
boundary divisor classes that generate
\(\Cl(\mathbb F_r)\).
\end{pauseandthink}

\subsection{Factorial and \texorpdfstring{\(\QQ\)}{Q}-factorial toric varieties}

\begin{definition}
A normal variety is \emph{factorial} if every Weil divisor is Cartier.

It is \emph{\(\Qfactorial\)} if every Weil divisor is \(\QCartier\);
that is, every Weil divisor has a positive multiple that is Cartier.
\end{definition}

These conditions ask whether arbitrary codimension-one
multiplicities admit integral or rational local equations.

\begin{theorem}[Cartier properties and the fan]
\label{thm:factorial-qfactorial-fan}
Let \(X_\Sigma\) be a normal toric variety.

\begin{enumerate}[label=(\alph*),wide=0pt]
\item The following are equivalent:
\[
X_\Sigma\text{ is smooth},
\qquad
X_\Sigma\text{ is factorial},
\qquad
\Pic(X_\Sigma)=\Cl(X_\Sigma).
\]

\item The following are equivalent:
\[
\Sigma\text{ is simplicial},
\qquad
X_\Sigma\text{ is }\Qfactorial,
\qquad
\Pic(X_\Sigma)
\text{ has finite index in }
\Cl(X_\Sigma).
\]
\end{enumerate}
\end{theorem}

\noindent
\textit{Reference.}
These are Cox--Little--Schenck, Propositions~4.2.6 and~4.2.7.

\begin{proof}[Conceptual proof sketch]
Work on an affine chart
\[
U_\sigma
\]
with ray generators
\[
u_1,\ldots,u_s.
\]
Prescribing a boundary divisor means prescribing arbitrary integers
\[
a_1,\ldots,a_s.
\]
The Cartier criterion asks whether the linear system
\[
\langle m,u_i\rangle=-a_i
\]
has a solution \(m\in M\).

It has an integral solution for every integer right-hand side exactly
when the \(u_i\) extend to a lattice basis. This is the smoothness
criterion.

It has a rational solution for every rational right-hand side exactly
when the \(u_i\) are linearly independent over \(\RR\). This says that
\(\sigma\) is simplicial.

Applying the local statement to every cone gives the theorem. The
statements involving \(\Pic\) and
\(\Cl\) merely rephrase whether all Weil classes are
Cartier or become Cartier after multiplication.
\end{proof}

\begin{observation}
Every strongly convex two-dimensional cone is simplicial. Therefore
every normal toric surface is \(\Qfactorial\), even though it may be
singular.

Thus, for toric surfaces,
\[
\text{smooth}
\Longleftrightarrow
\text{factorial},
\]
but \(\Qfactorial\) is automatic.
\end{observation}

The cyclic quotient
\[
X_n=\AAA^2/\tfrac1n(1,n-1)
\]
is a good test case. For \(n>1\), it is singular and therefore not
factorial. Its cone is simplicial, so it is nevertheless
\(\Qfactorial\).

\subsection{Gorenstein and \texorpdfstring{\(\QQ\)}{Q}-Gorenstein cones}

\begin{definition}
A normal variety is \emph{Gorenstein} if its canonical divisor is
Cartier. It is \emph{\(\QGorenstein\)} if its canonical divisor is
\(\QCartier\).
\end{definition}

Because
\[
K_{X_\Sigma}
=
-\sum_\rho D_\rho,
\]
its coefficient along every boundary divisor is \(-1\). The Cartier
criterion therefore becomes a height-one condition.

\begin{proposition}[The height-one criterion]
\label{prop:gorenstein-height-one}
Let \(X_\Sigma\) be a normal toric variety.

\begin{enumerate}[label=(\alph*),wide=0pt]
\item \(X_\Sigma\) is Gorenstein if and only if, for every maximal
cone \(\sigma\), there exists \(m_\sigma\in M\) such that
\[
\langle m_\sigma,u_\rho\rangle=1
\qquad
\text{for every }\rho\in\sigma(1).
\]

\item \(X_\Sigma\) is \(\QGorenstein\) if and only if the same
condition holds with
\[
m_\sigma\in M_\QQ.
\]
\end{enumerate}
\end{proposition}

\noindent
\textit{Reference.}
This is Cox--Little--Schenck, Proposition~8.2.12; the Gorenstein
terminology is Definition~8.2.14.

The condition says that the primitive ray generators of each maximal
cone lie on an affine hyperplane
\[
\{u\in N_\RR:\langle m_\sigma,u\rangle=1\}.
\]
The normal vector must be integral for Gorensteinness and merely
rational for \(\QGorenstein\).

\begin{figure}[ht]
\centering
\begin{tikzpicture}[scale=1.02]
  % A_{n-1} cone
  \begin{scope}[xshift=-4.5cm]
    \fill[DictionaryBlue!9]
      (0,0)--(3.0,0)--(0,3.0)--cycle;

    \draw[very thick,-{Latex[length=2.3mm]}]
      (0,0)--(3.25,0);
    \draw[very thick,-{Latex[length=2.3mm]}]
      (0,0)--(0,3.25);

    \draw[very thick,PredictionTeal]
      (0,3.0)--(3.0,0);

    \foreach \i in {0,...,5} {
      \fill[PredictionTeal]
        ({3*\i/5},{3*(5-\i)/5}) circle (2.2pt);
    }

    \node[left] at (0,3.0) {\(u_0=e_2\)};
    \node[below] at (3.0,0) {\(u_n=e_1\)};
    \node[PredictionTeal,above right] at (1.65,1.65)
      {\(x+y=1\)};
    \node at (1.5,0.62)
      {\(\langle(1,1),u_i\rangle=1\)};
    \node[font=\bfseries] at (1.4,3.75)
      {\(A_{n-1}\): integral height one};
  \end{scope}

  % Weighted projective singular cone
  \begin{scope}[xshift=4.6cm]
    \coordinate (O) at (0,0);
    \coordinate (A) at (3.0,0);
    \coordinate (B) at (-1.5,3.0);

    \fill[ExampleGold!10] (O)--(A)--(B)--cycle;
    \draw[very thick,-{Latex[length=2.3mm]}] (O)--(3.25,0);
    \draw[very thick,-{Latex[length=2.3mm]}] (O)--(-1.65,3.3);
    \draw[very thick,WarningRed] (A)--(B);

    \fill[WarningRed] (A) circle (2.2pt);
    \fill[WarningRed] (B) circle (2.2pt);

    \node[below] at (A) {\((1,0)\)};
    \node[above left] at (B) {\((-1,n)\)};
    \node[WarningRed,rotate=-34] at (0.75,1.82)
      {\(x+\frac{2}{n}y=1\)};
    \node[align=center] at (0.65,0.5)
      {integral normal only when\\ \(n\mid2\)};
    \node[font=\bfseries] at (0.75,3.75)
      {\(\PP(1,1,n)\): rational height one};
  \end{scope}
\end{tikzpicture}
\caption{Gorensteinness is an integral height-one condition. For the
\(A_{n-1}\) cone, the character \((1,1)\) is integral in the intrinsic
character lattice. For the singular cone of \(\PP(1,1,n)\), the
height-one functional is \((1,2/n)\), which is integral only for
\(n=1,2\).}
\label{fig:gorenstein-height-one}
\end{figure}

\subsection{Two singular comparisons}

\begin{example}[The \(A_{n-1}\) singularity is Gorenstein]
For
\[
X_n=\AAA^2/\tfrac1n(1,n-1),
\]
the character lattice is
\[
M
=
\{(a,b)\in\ZZ^2:a-b\equiv0\pmod n\}.
\]
The element
\[
m=(1,1)
\]
belongs to \(M\), and
\[
\langle m,e_1\rangle
=
\langle m,e_2\rangle
=
1.
\]
Hence \(X_n\) is Gorenstein.

This recovers geometrically the special role of the equation
\[
uv-w^n=0.
\]
The singularity is not smooth for \(n>1\), but its canonical divisor
is Cartier.

Moreover, the rays in the resolution from the Chapter~3 challenge,
\[
u_i=
\left(\frac{i}{n},\frac{n-i}{n}\right),
\]
all satisfy
\[
\langle(1,1),u_i\rangle=1.
\]
The new rays remain on the same height-one line.
\end{example}

\begin{observation}
For the general quotient
\[
\AAA^2/\tfrac1n(1,q),
\]
the unique rational functional taking value \(1\) on \(e_1,e_2\) is
\[
m=(1,1).
\]
It lies in the intrinsic character lattice exactly when
\[
1+q\equiv0\pmod n.
\]
For \(1\leq q<n\), this means
\[
q=n-1.
\]

Thus the Gorenstein quotients in this family are precisely the
\(A_{n-1}\) singularities. Every member of the family is
\(\QGorenstein\), and its Gorenstein index is
\[
\frac{n}{\gcd(n,q+1)}.
\]
\end{observation}

\begin{example}[The weighted projective plane]
The fan of \(\PP(1,1,n)\) has a singular cone generated by
\[
u_0=(1,0),
\qquad
u_1=(-1,n).
\]
The height-one equations force
\[
m=\left(1,\frac{2}{n}\right).
\]
Therefore
\[
\PP(1,1,n)
\text{ is Gorenstein}
\quad\Longleftrightarrow\quad
n\mid2.
\]
Hence the singular surface
\[
\PP(1,1,2)
\]
is Gorenstein, while \(\PP(1,1,n)\) is only \(\QGorenstein\) for
\(n>2\).

The fan is simplicial for every \(n\), so all these weighted
projective planes are \(\Qfactorial\).
\end{example}

\subsection{A recognition table}

\begin{center}
\small
\renewcommand{\arraystretch}{1.3}
\begin{tabularx}{\linewidth}{@{}
  >{\raggedright\arraybackslash}p{0.22\linewidth}
  >{\raggedright\arraybackslash}p{0.31\linewidth}
  >{\raggedright\arraybackslash}X@{}}
\toprule
\textbf{Property}
&
\textbf{Divisor formulation}
&
\textbf{Fan formulation}
\\
\midrule
smooth / factorial
&
every Weil divisor is Cartier
&
every cone is smooth
\\
\(\Qfactorial\)
&
every Weil divisor is \(\QCartier\)
&
every cone is simplicial
\\
Gorenstein
&
\(K_X\) is Cartier
&
each maximal cone has an integral height-one functional
\\
\(\QGorenstein\)
&
\(K_X\) is \(\QCartier\)
&
each maximal cone has a rational height-one functional
\\
\bottomrule
\end{tabularx}
\end{center}

The implications should not be confused:
\[
\text{smooth}
\Longrightarrow
\begin{cases}
\text{factorial},\\
\text{Gorenstein},\\
\Qfactorial.
\end{cases}
\]
But a singular toric variety may still be Gorenstein, \(\Qfactorial\),
or both.

The chapter has translated divisors into coefficients on rays,
relations from characters, local linear equations, and support
functions. The next step is to ask which characters become global
sections.

\begin{guidingquestion}
Given
\[
D=\sum_\rho a_\rho D_\rho,
\]
why should the inequalities
\[
\langle m,u_\rho\rangle\geq-a_\rho
\qquad
\text{for every }\rho
\]
select exactly the characters \(\character{m}\) for which
\[
\operatorname{div}(\character{m})+D\geq0?
\]

What geometric object in \(M_\RR\) is cut out by these inequalities?
\end{guidingquestion}

\clearpage
\section*{What this chapter has accomplished}
\addcontentsline{toc}{section}{What this chapter has accomplished}
\markboth{What this chapter has accomplished}{What this chapter has accomplished}

The rays of a fan record the irreducible components of the toric
boundary:
\[
\rho\in\Sigma(1)
\quad\longleftrightarrow\quad
D_\rho=\overline{O(\rho)}.
\]
Characters then measure zeros and poles along those components:
\[
\operatorname{div}(\character{m})
=
\sum_{\rho\in\Sigma(1)}
\langle m,u_\rho\rangle D_\rho.
\]

From this formula, the chapter developed the sequence
\[
\boxed{
\begin{aligned}
\text{rays}
&\longrightarrow
\text{invariant prime divisors},
\\
\text{characters}
&\longrightarrow
\text{principal relations},
\\
\text{Cartier data}
&\longrightarrow
\text{compatible local equations},
\\
\text{support functions}
&\longrightarrow
\text{one piecewise-linear encoding}.
\end{aligned}
}
\]

The two basic divisor groups are organized by
\[
M
\longrightarrow
\Div_T(X_\Sigma)
\longrightarrow
\Cl(X_\Sigma)
\longrightarrow0
\]
and
\[
M
\longrightarrow
\operatorname{SF}(\Sigma,N)
\longrightarrow
\Pic(X_\Sigma)
\longrightarrow0.
\]
The difference between \(\Cl\) and
\(\Pic\) measures the failure of Weil divisors to admit
compatible local equations.

Finally,
\[
K_{X_\Sigma}
=
-\sum_{\rho\in\Sigma(1)}D_\rho
\]
turned singularity properties into fan conditions:
\[
\begin{array}{c@{\quad\longleftrightarrow\quad}c}
\text{factorial}
&
\text{smooth fan},
\\
\Qfactorial
&
\text{simplicial fan},
\\
\text{Gorenstein}
&
\text{integral height-one data},
\\
\QGorenstein
&
\text{rational height-one data}.
\end{array}
\]

The next chapter begins with a divisor
\[
D=\sum_\rho a_\rho D_\rho
\]
and asks which characters satisfy
\[
\operatorname{div}(\character{m})+D\geq0.
\]
The answer will be encoded by lattice points in the polyhedron
\[
P_D
=
\left\{
m\in M_\RR:
\langle m,u_\rho\rangle\geq-a_\rho
\text{ for every }\rho
\right\}.
\]

\begin{exercises}[Comprehension exercises]
\markboth{Comprehension exercises}{Comprehension exercises}

\item Let \(X=\PP^1\), with boundary divisors
\[
D_0=\{0\},
\qquad
D_\infty=\{\infty\}.
\]
\begin{enumerate}[label=(\alph*),wide=0pt]
\item Use the character \(t\) to compute
\[
\Cl(\PP^1).
\]
\item Write Cartier data for \(D_0\) on the two affine charts.
\item Determine the support function of \(D_0\).
\item Recover
\[
K_{\PP^1}\sim-2[D_0].
\]
\end{enumerate}

\item Let
\[
X=\AAA^1\times\Gm
\]
with its usual toric structure.
\begin{enumerate}[label=(\alph*),wide=0pt]
\item Describe its fan and its invariant prime divisors.
\item Compute the map
\[
M\longrightarrow\Div_T(X).
\]
\item Determine \(\Cl(X)\).
\item Find the kernel of the divisor map and explain it using the
torus factor.
\end{enumerate}

\item For the fan of \(\PP^2\), let
\[
D=a_0D_0+a_1D_1+a_2D_2.
\]
\begin{enumerate}[label=(\alph*),wide=0pt]
\item Show directly that every such invariant Weil divisor is Cartier.
\item Find Cartier data on the three maximal cones.
\item Prove
\[
[D]=(a_0+a_1+a_2)H.
\]
\item Characterize the principal invariant divisors among these
triples.
\end{enumerate}

\item For \(\PP^1\times\PP^1\), order the invariant boundary divisors
cyclically as
\[
D_1,D_2,D_3,D_4.
\]
\begin{enumerate}[label=(\alph*),wide=0pt]
\item Recover the two relations
\[
D_1\sim D_3,
\qquad
D_2\sim D_4.
\]
\item Compute the support function of \(D_1+D_2\).
\item Express the canonical class in the basis
\[
[D_1],[D_2].
\]
\end{enumerate}

\item Let \(\mathbb F_r\) have rays
\[
(1,0),\quad
(0,1),\quad
(-1,r),\quad
(0,-1),
\]
with corresponding divisors
\[
D_1,D_2,D_3,D_4.
\]
\begin{enumerate}[label=(\alph*),wide=0pt]
\item Compute the two principal relations among the \(D_i\).
\item Determine \(\Cl(\mathbb F_r)\).
\item Express \(K_{\mathbb F_r}\) using two chosen generators.
\item Check that every invariant Weil divisor is Cartier.
\end{enumerate}

\item Consider the affine cone
\[
\sigma=\Cone{(1,0),(1,d)}
\subseteq\RR^2,
\qquad
d\geq2,
\]
in the lattice \(\ZZ^2\).
\begin{enumerate}[label=(\alph*),wide=0pt]
\item Compute \(\Cl(U_\sigma)\).
\item Show that each invariant prime divisor is \(\QCartier\).
\item Determine the smallest positive multiple of the divisor
corresponding to \((1,0)\) that is Cartier.
\item Decide when \(U_\sigma\) is Gorenstein.
\end{enumerate}

\item Assign integers
\[
c_0,c_1,c_2
\]
to the three rays of the fan of \(\PP^2\).
\begin{enumerate}[label=(\alph*),wide=0pt]
\item Show that there is a unique support function taking those values
on the rays.
\item Write its three linear pieces explicitly.
\item Which choice of \(c_i\) gives a global linear function?
\item Translate your answer into a statement about principal
divisors.
\end{enumerate}

\item Let \(\Sigma\) consist of the cone
\[
\Cone{e_1,e_2,e_3,e_1+e_2-e_3}
\subseteq\RR^3
\]
and all its faces.
\begin{enumerate}[label=(\alph*),wide=0pt]
\item Explain why the maximal cone is not simplicial.
\item Construct a Weil divisor whose coefficients on the four rays
cannot arise from a rational linear function.
\item Deduce that \(X_\Sigma\) is not \(\Qfactorial\).
\end{enumerate}

\item Let \(X=\PP(1,1,n)\).
\begin{enumerate}[label=(\alph*),wide=0pt]
\item Use the fan to verify that \(X\) is \(\Qfactorial\).
\item Compute the rational height-one functional on its singular cone.
\item Determine the Cartier index of \(K_X\).
\item Verify that \(X\) is Gorenstein precisely when \(n\mid2\).
\end{enumerate}

\item Let \(\Sigma'\) be a refinement of a fan \(\Sigma\), with the
same support, and let
\[
\pi:X_{\Sigma'}\longrightarrow X_\Sigma
\]
be the induced proper birational toric morphism.
Suppose \(D\) is a torus-invariant Cartier divisor on \(X_\Sigma\)
with support function \(\psi_D\).
\begin{enumerate}[label=(\alph*),wide=0pt]
\item Explain why the same function is a support function on
\(\Sigma'\).
\item Show that the corresponding divisor on \(X_{\Sigma'}\) is
\(\pi^*D\).
\item If \(\rho'\) is a new ray, determine its coefficient in
\(\pi^*D\) from \(\psi_D(u_{\rho'})\).
\end{enumerate}

\end{exercises}

\clearpage
\section*{Challenge exercise}
\addcontentsline{toc}{section}{Challenge exercise}
\markboth{Challenge exercise}{Challenge exercise}

\noindent
\textbf{Divisors and the crepant resolution of an
\(\boldsymbol{A_{n-1}}\) singularity.}

\medskip

Assume \(\Bbbk=\CC\), and consider
\[
X_n
=
\AAA^2/\tfrac1n(1,n-1)
\simeq
V(uv-w^n)\subseteq\AAA^3.
\]
Its fan consists of the first-quadrant cone
\[
\sigma=\Cone{e_2,e_1}
\]
in
\[
N
=
\ZZ^2+
\ZZ\left(\frac1n(1,n-1)\right).
\]
The intrinsic character lattice is
\[
M
=
\left\{
(a,b)\in\ZZ^2:
a-b\equiv0\pmod n
\right\}.
\]

As in the Chapter~3 challenge, define
\[
u_i
=
\left(\frac{i}{n},\frac{n-i}{n}\right),
\qquad
0\leq i\leq n.
\]
Thus
\[
u_0=e_2,
\qquad
u_n=e_1.
\]
Let \(\Sigma'\) be the smooth fan with maximal cones
\[
\sigma_i=\Cone{u_{i-1},u_i},
\qquad
1\leq i\leq n,
\]
and let
\[
\pi:Y=X_{\Sigma'}\longrightarrow X_n
\]
be the resulting resolution.

Write
\[
D_0=V(\Cone{u_0}),
\qquad
D_n=V(\Cone{u_n}),
\]
for the strict transforms of the two original boundary divisors, and
\[
E_i=V(\Cone{u_i}),
\qquad
1\leq i\leq n-1,
\]
for the exceptional curves.

The goal is to understand how divisor classes and canonical divisors
behave under the resolution.

\clearpage
\begin{figure}[ht]
\centering
\begin{tikzpicture}[scale=0.92]
  % singular cone
  \begin{scope}[xshift=-4.7cm]
    \fill[WarningRed!9]
      (0,0)--(3.0,0)--(0,3.0)--cycle;
    \draw[very thick,-{Latex[length=2.3mm]}]
      (0,0)--(3.25,0);
    \draw[very thick,-{Latex[length=2.3mm]}]
      (0,0)--(0,3.25);
    \draw[very thick,PredictionTeal]
      (0,3.0)--(3.0,0);

    \node[left] at (0,3.0) {\(u_0=e_2\)};
    \node[below] at (3.0,0) {\(u_n=e_1\)};
    \node[PredictionTeal,above right] at (1.55,1.55)
      {\(\langle(1,1),u\rangle=1\)};
    \node[font=\bfseries] at (1.45,3.75)
      {the singular cone};
  \end{scope}

  \draw[very thick,-{Latex[length=2.8mm]}]
    (-0.75,0.45)--(0.75,0.45);
  \node[above] at (0,0.45) {refine};

  % resolution fan and chain
  \begin{scope}[xshift=4.6cm]
    \draw[very thick,-{Latex[length=2.2mm]}]
      (0,0)--(3.15,0);
    \draw[very thick,-{Latex[length=2.2mm]}]
      (0,0)--(0,3.15);

    \foreach \i in {1,...,4} {
      \draw[very thick,PredictionTeal,-{Latex[length=2.2mm]}]
        (0,0)--({3*\i/5},{3*(5-\i)/5});
    }
    \draw[very thick,PredictionTeal]
      (0,3.0)--(3.0,0);

    \node[left] at (0,3.0) {\(u_0\)};
    \node[above left] at (0.6,2.4) {\(u_1\)};
    \node[above] at (1.2,1.8) {\(u_2\)};
    \node at (1.8,1.25) {\(\cdots\)};
    \node[above right] at (2.4,0.6) {\(u_{n-1}\)};
    \node[below] at (3.0,0) {\(u_n\)};

    \node[font=\bfseries] at (1.5,3.75)
      {all rays remain at height one};

    \node[
      draw=PredictionTeal,
      rounded corners,
      inner sep=2pt
    ] (E1) at (0.15,-1.0) {\(E_1\)};
    \node[
      draw=PredictionTeal,
      rounded corners,
      inner sep=2pt
    ] (E2) at (0.95,-1.0) {\(E_2\)};
    \node at (1.75,-1.0) {\(\cdots\)};
    \node[
      draw=PredictionTeal,
      rounded corners,
      inner sep=2pt
    ] (En) at (2.75,-1.0) {\(E_{n-1}\)};
    \draw[thick] (E1)--(E2);
    \draw[thick] (E2)--(1.48,-1.0);
    \draw[thick] (2.05,-1.0)--(En);
  \end{scope}
\end{tikzpicture}
\caption{The resolution inserts the intermediate lattice rays. Every
ray lies on the same integral height-one line, which will force the
resolution to preserve the canonical divisor.}
\label{fig:challenge-crepant-an-resolution}
\end{figure}

\begin{enumerate}[label=\textbf{Hint \arabic*.},wide=0pt]

\item \textbf{The class group of the singularity.}

Let
\[
B_0=V(\Cone{u_0}),
\qquad
B_n=V(\Cone{u_n})
\]
be the invariant boundary divisors on \(X_n\). Show that
\[
m_0=(1,1),
\qquad
m_1=(n,0)
\]
form a basis of \(M\), and compute
\[
\operatorname{div}_{X_n}(\character{m_0})
=
B_0+B_n,
\]
\[
\operatorname{div}_{X_n}(\character{m_1})
=
nB_n.
\]
Deduce that
\[
\Cl(X_n)
\simeq
\ZZ/n\ZZ,
\]
generated by \([B_n]\), with
\[
[B_0]=-[B_n].
\]

\item \textbf{The class group upstairs.}

On \(Y\), write
\[
D_i=
\begin{cases}
D_0,&i=0,\\
E_i,&1\leq i\leq n-1,\\
D_n,&i=n.
\end{cases}
\]
Verify
\[
\langle m_0,u_i\rangle=1,
\qquad
\langle m_1,u_i\rangle=i.
\]
Hence
\[
\operatorname{div}_Y(\character{m_0})
=
\sum_{i=0}^{n}D_i,
\]
and
\[
\operatorname{div}_Y(\character{m_1})
=
\sum_{i=0}^{n}iD_i.
\]

Use these two relations to show that
\[
\Cl(Y)
\simeq
\ZZ^{n-1}.
\]
Why is the answer torsion-free without performing a Smith normal form
calculation?

\item \textbf{What the exceptional curves generate.}

Let \(L\subseteq\Cl(Y)\) be the subgroup generated by
\[
[E_1],\ldots,[E_{n-1}].
\]
Show that
\[
\Cl(Y)/L
\simeq
\ZZ/n\ZZ.
\]
Compare this quotient with \(\Cl(X_n)\).

Explain informally why contracting the exceptional chain should
recover the torsion class group of the singular surface.

\item \textbf{Cartier data for one exceptional curve.}

Fix \(1\leq i\leq n-1\). Since \(Y\) is smooth, \(E_i\) is Cartier.
Show that the only maximal cones on which a nontrivial local equation
is needed are
\[
\sigma_i=\Cone{u_{i-1},u_i},
\qquad
\sigma_{i+1}=\Cone{u_i,u_{i+1}}.
\]

Define
\[
m_i^-=(i-1-n,i-1),
\qquad
m_i^+=(n-i-1,-i-1).
\]
Verify that both belong to \(M\) and that
\[
\begin{array}{c|cc}
&
u_{i-1}
&
u_i
\\
\hline
m_i^-
&
0
&
-1
\end{array},
\qquad
\begin{array}{c|cc}
&
u_i
&
u_{i+1}
\\
\hline
m_i^+
&
-1
&
0.
\end{array}
\]
Conclude that Cartier data for \(E_i\) are
\[
m_{\sigma_i}=m_i^-,
\qquad
m_{\sigma_{i+1}}=m_i^+,
\qquad
m_\sigma=0
\]
on every other maximal cone.

Write the corresponding local character equations and check their
compatibility on overlaps.

\item \textbf{Total transforms of the boundary.}

Since
\[
nB_n
=
\operatorname{div}_{X_n}(\character{m_1}),
\]
pulling back this principal divisor gives
\[
\pi^*(nB_n)
=
\operatorname{div}_Y(\character{m_1}).
\]
Deduce
\[
\boxed{
\pi^*(nB_n)
=
nD_n+\sum_{i=1}^{n-1}iE_i.
}
\]

Similarly, use
\[
(0,n)=nm_0-m_1
\]
to show
\[
\boxed{
\pi^*(nB_0)
=
nD_0+\sum_{i=1}^{n-1}(n-i)E_i.
}
\]

How do the increasing and decreasing coefficients reflect the
positions of the rays in the fan?

\item \textbf{The canonical divisor downstairs.}

Use the toric canonical-divisor formula to write
\[
K_{X_n}=-B_0-B_n.
\]
Show that
\[
K_{X_n}
=
\operatorname{div}_{X_n}(\character{-m_0}).
\]
Thus the canonical divisor is not merely Cartier: this chosen
representative is principal.

Relate this calculation to the fact that the two primitive ray
generators lie on the integral height-one line
\[
\langle m_0,u\rangle=1.
\]

\item \textbf{The canonical divisor upstairs.}

Since every ray of \(\Sigma'\) satisfies
\[
\langle m_0,u_i\rangle=1,
\]
show that
\[
K_Y
=
-\sum_{i=0}^{n}D_i
=
\operatorname{div}_Y(\character{-m_0}).
\]
Now use pullback of principal divisors to conclude
\[
\boxed{
K_Y=\pi^*K_{X_n}.
}
\]

A resolution satisfying this equality is called \emph{crepant}.
Explain, using only the fan picture, why no exceptional divisor
appears with a nonzero correction coefficient.

\item \textbf{The support-function interpretation.}

The support function of \(K_{X_n}\) is the global linear function
\[
\psi_K(u)=\langle m_0,u\rangle.
\]
Regard the same function on the refined fan \(\Sigma'\).

\begin{enumerate}[label=(\alph*),wide=0pt]
\item What coefficient does the pullback divisor assign to each new
ray \(u_i\)?
\item Why does this agree with the coefficient of \(E_i\) in \(K_Y\)?
\item Explain why the equality
\[
K_Y=\pi^*K_{X_n}
\]
is equivalent to saying that every inserted ray remains at height
one.
\end{enumerate}

\item \textbf{A concrete case.}

Work out the case \(n=5\).

\begin{enumerate}[label=(\alph*),wide=0pt]
\item Draw the six rays \(u_0,\ldots,u_5\).
\item Write the two principal relations in
\(\Div_T(Y)\).
\item Compute Cartier data for \(E_2\).
\item Write the total transforms of \(5B_0\) and \(5B_5\).
\item Verify directly that
\[
K_Y=\pi^*K_{X_5}.
\]
\end{enumerate}

\end{enumerate}

\begin{tcolorbox}[
  colback=PredictionTeal!5,
  colframe=PredictionTeal,
  title={Beyond the \(A_{n-1}\) case},
  fonttitle=\bfseries
]
For a general cyclic quotient
\(\AAA^2/\tfrac1n(1,q)\), the canonical support function is rational
and takes value \(1\) on the two original rays. A new primitive ray
\(u\) contributes
\[
\langle m_K,u\rangle-1
\]
to the coefficient of its exceptional curve in
\(K_Y-\pi^*K_X\).

For \(A_{n-1}\), every inserted ray remains at height \(1\), so all
these corrections vanish.
\end{tcolorbox}

\noindent
\textit{Reference.}
The divisor class sequence is Cox--Little--Schenck,
Theorem~4.1.3, and support functions are developed in \S4.2.
Canonical divisors and the height-one criterion are treated in
\S8.2. The toric resolution of cyclic quotient surface singularities
and its exceptional chain are developed in Chapter~10, especially
\S\S10.2 and~10.4.

\clearpage
\section*{Notes and references}
\addcontentsline{toc}{section}{Notes and references}
\markboth{Notes and references}{Notes and references}

The invariant divisor sequence, the Cartier criterion, and support
functions follow Cox--Little--Schenck, Chapter~4. The formula
\[
K_{X_\Sigma}
=
-\sum_{\rho\in\Sigma(1)}D_\rho
\]
and the Gorenstein height-one criterion are developed in
Cox--Little--Schenck, \S8.2.

The chapter deliberately stops before global sections and positivity.
Those topics require one further translation:
\[
\operatorname{div}(\character{m})+D\geq0
\quad\Longleftrightarrow\quad
\langle m,u_\rho\rangle\geq-a_\rho
\text{ for every }\rho.
\]
Chapter~5 will turn these inequalities into polyhedra, interpret their
lattice points as sections, and connect convex support functions to
projective embeddings.

The challenge exercise returns to the \(A_{n-1}\) resolution from
Chapter~3. It illustrates the principle that a toric resolution is
crepant precisely when the new rays remain on the canonical
height-one hyperplane.
